\subsection{里斯自同态的扰动}

引理 1. --- 设 E 是巴拿赫空间。设 p 和 q 是 E 上的连续投影，满足 \(\|q - p\| < 1\) 。则 p 诱导从 \(\operatorname{Im}(q)\) 到 \(\operatorname{Im}(p)\) 上的一个同构，而 \(1_{E} - p\) 诱导从 \(\operatorname{Ker}(q)\) 到 \(\operatorname{Ker}(p)\) 上的一个同构。

考察连续线性映射 \(u\colon x\mapsto p(x)\) （从 \(\operatorname{Im}(q)\) 到 \(\operatorname{Im}(p)\) 中）与 \(v\colon y\mapsto q(y)\) （从 \(\operatorname{Im}(p)\) 到 \(\operatorname{Im}(q)\) 中）。对每个 \(x\in\operatorname{Im}(q)\) ，有 \(x=q(x)\) ，由此

\[
\begin{array}{c} (q - p) ^ {2} (x) = q ^ {2} (x) - q (p (x)) - p (q (x)) + p ^ {2} (x) \\ = x - q (p (x)) = x - v (u (x)). \end{array}
\]

由此可得 \(\|1_{E}-v\circ u\|\leqslant\|q-p\|^{2}<1\) 。由 I, p. 22 的推论 1，\(v\circ u\) 是 \(\operatorname{Im}(q)\) 的一个自同构。类似地可证 \(u\circ v\) 是 \(\operatorname{Im}(p)\) 的一个自同构。这蕴含 u 是从 \(\operatorname{Im}(q)\) 到 \(\operatorname{Im}(p)\) 上的一个同构，从而得出引理的第一个断言。第二个断言由此通过把 p 和 q 分别换成 \(1_{E}-p\) 和 \(1_{E}-q\) 而得到。

定理 2. --- 设 E 是巴拿赫空间。E 的里斯自同态所成的集 \(\mathcal{R}(\mathrm{E})\) 在 \(\mathcal{L}(\mathrm{E})\) 中是开的。把 \(\mathcal{R}(\mathrm{E})\) 的元素映为其幂零空间的维数的映射在 \(\mathcal{L}(\mathrm{E})\) 上是上半连续的。

定理由下面更精确的命题得出：

命题 5. --- 设 E 是巴拿赫空间，\(u_{0}\) 是 E 的一个里斯自同态。设 N（相应地，I）是 \(u_{0}\) 的幂零空间（相应地，余幂零空间）。用 d 表示 N 的维数。存在 \(u_{0}\) 在 \(\mathcal{L}(\mathrm{E})\) 中的一个开邻域 U 和一个解析映射 \(\pi\colon\mathrm{U}\to\mathcal{L}(\mathrm{E})\) ，使得

\begin{enumerate}
\def\labelenumi{(\roman{enumi})}
\item
  自同态 \(\pi(u_0)\) 是以 N 为像、以 I 为核的投影；
\item
  对每个 \(u \in U\) ，线性映射 \(\pi(u)\) 是秩为 d 的投影，它属于 u 的双交换子，特别地与 u 交换。此外，u 诱导 \(\operatorname{Ker}(\pi(u))\) 的一个自同构；
\item
  U 的每个元素都是里斯自同态，其幂零空间的维数 \(\leqslant d\) 。
\end{enumerate}

若 K = C，用 \(\mathrm{Sp}(u_{0})\) 表示 \(u_{0}\) 关于代数 \(\mathcal{L}(\mathrm{E})\) 的谱。当 \(K = R\) 时，用 \(\mathrm{Sp}(u_{0})\) 表示复谱 \(\mathrm{Sp}_{\mathcal{L}(\mathrm{E}_{(\mathbf{C})})}((u_{0})_{(\mathbf{C})})\) ，即 \(u_{0}\) 的复谱（I, p. 85, 第 13 目）。由 III, p. 54 的命题 14，点 0 在 \(\{0\} \cup \mathrm{Sp}(u_{0})\) 中是孤立的。设 r \textgreater{} 0 是实数，使得 \(\mathrm{Sp}(u_{0}) - \{0\}\) 的每个元素的模都 \textgreater{} r。用 V 表示 C 中由绝对值 ≠ r 的复数所组成的开集；用 f 表示 V 上如下定义的全纯函数：\(f(z) = 0\) （若 \(|z| > r\) ），\(f(z) = 1\) （若 \(|z| < r\) ）。若 K = C（相应地，K = R），用 U' 表示 \(\mathcal{L}(\mathrm{E})\) 中谱（相应地，复谱）包含于 V 的元素 u 所成的集。集 U' 是 \(u_{0}\) 在 \(\mathcal{L}(\mathrm{E})\) 中的一个开邻域，且映射 \(u \mapsto f(u)\) （从 U' 到 \(\mathcal{L}(\mathrm{E})\) 中）是全纯的（I, p. 76, 命题 10 及 I, p. 85, 第 13 目）。

设 \(u \in \mathrm{U}'\) 。自同态 \(f(u)\) 是谱投影 \(e_0(u)\) ；它与 \(u\) 交换，且 \(u\) 通过过渡到子空间诱导 \(\operatorname{Ker}(e_0(u))\) 的一个自同构（参见 III, p. 53, 第 6 目）。

投影 \(e_{0}(u_{0})\) 以 N 为像、以 I 为核（III, p. 54, 命题 14）；它的秩是 d。由引理 1，由 \(u \in U'\) 中使 \(e_{0}(u)\) 的秩为 d 的元素所成的集 U 是 \(u_{0}\) 的一个开邻域。集 U 与映射 \(\pi: u \mapsto e_{0}(u)\) （从 U 到 \(\mathcal{L}(\mathrm{E})\) 中）满足命题的条件 (i) 与 (ii)。

设 \(u \in U\) 。由于 u 诱导 \(\operatorname{Ker}(e_{0}(u))\) （它是 E 的一个闭的有限余维数向量子空间）的一个自同构，自同态 u 是 E 的里斯自同态（III, p. 48, 命题 8）；u 的幂零空间包含于 \(\pi(u) = e_{0}(u)\) 的像中，且其维数 \(\leqslant d\) 。

\subsection{连续线性映射的余范数}

设 E 和 F 是赋范空间，\(u: E \rightarrow F\) 是连续线性映射，N 是 u 的核，I 是它的像。用 p 表示从 E 到 E/N 上的典范满射，用 i 表示从 I 到 F 中的典范单射。用 \(\widetilde{u}\) 表示从 E/N 到 I 上的双射线性映射，它满足 \(u = i \circ \widetilde{u} \circ p\) 。向量空间 E/N 赋有商范数，即

\[
\| y \| = \inf _ {x \in p^{-1}(y)} \| x \|\tag{2}
\]

对一切 \(y \in E/N\) 成立。映射 \(\widetilde{u}\) 连续且 \(\|u\| = \|\widetilde{u}\|\) ，由此

\[
\| u\| = \sup_{\substack{y\in \mathrm{E / N}\\ y\neq 0}}\frac{\|\widetilde{u} (y)\|}{\|y\|},\tag{3}
\]

其中上确界在 \(\overline{R}_{+}\) 中取。u 的余范数是指

\[
((u)) = \inf_{\substack{y\in \mathrm{E} / \mathrm{N}\\ y\neq 0}}\frac{\| \widetilde{u}(y)\|}{\|y\|},\tag{4}
\]

其中下确界在 \(\overline{R}_{+}\) 中取。因此我们有

\[
((u)) \| y \| \leqslant \| \widetilde {u} (y) \| \leqslant \| u \| \| y \|\tag{5}
\]

对 E/N 的每个元素 \(y\) 成立。令 \(v = i \circ \widetilde{u}\) 。我们有 \(u = v \circ p\) ，并且

\[
((u)) = \inf_{\substack{y\in \mathrm{E / N}\\ y\neq 0}}\frac{\|v(y)\|}{\|y\|} = \inf_{\substack{y\in \mathrm{E / N}\\ \|y\| = 1}}\| v(y)\| .\tag{6}
\]

当 \(u\) 是零映射时，空间 E/N 缩减为 0，且有 \(((u)) = +\infty\) 与 \(\| u \| = 0\) 。当 \(u \neq 0\) 时，由 (3) 与 (4) 推出不等式

\[
0 \leqslant ((u)) \leqslant \| u \| <   + \infty .\tag{7}
\]

当 u 单射时，我们有

\[
((u)) = \inf_{\substack{x\in \mathrm{E}\\ x\neq 0}}\frac{\|u(x)\|}{\|x\|},\tag{8}
\]

且对一切 \(x \in \mathrm{E}\) ，我们有

\[
((u)) \| x \| \leqslant \| u (x) \| \leqslant \| u \| \| x \|.\tag{9}
\]

用 \(j\) 表示赋范空间 F 到其完备化 \(\widehat{\mathbf{F}}\) 中的典范单射。我们有 \(((u)) = ((j \circ u))\) 。

注记. --- 按定义，\(((u)) > 0\) 等价于双射线性映射 \(\widetilde{u}\) 是双连续的，亦即等价于 u 是严格态射（TG, III, p. 16）。此时我们有

\[
((u)) = \| \widetilde {u} ^ {- 1} \| ^ {- 1}.\tag{10}
\]

引理 2. --- 设 c 是实数。若 c \textless{} ((u))，则对每个元素 \(z \in \operatorname{Im}(u)\) ，存在 E 的元素 x，使得 \(u(x) = z\) 且 \(c \|x\| \leqslant \|z\|\) 。反之，若对每个 \(z \in \operatorname{Im}(u)\) 存在 \(x \in E\) ，使得 \(u(x) = z\) 且 \(c \|x\| \leqslant \|z\|\) ，则 \(c \leqslant ((u))\) 。

这是公式 (2)、公式 (5) 以及 \(u\) 的余范数定义的推论。

命题 6. --- 设 E 和 F 是赋范空间，\(u \in \mathcal{L}(\mathrm{E};\mathrm{F})\) 。用 B 表示 E 中范数 \(< 1\) 的元素所成的集。令

\[
\mathrm{P} = u (\mathrm{E}) - u (\mathrm{B}), \quad \mathrm{Q} = u (\mathrm{E}) - (\overline {{u (\mathrm{B})}} \cap u (\mathrm{E})).
\]

u 的余范数等于在 F 中从 0 到 P 的距离。若赋范空间 E 完备或 u 是严格态射，则 u 的余范数等于在 F 中从 0 到 Q 的距离。

设 N 是 u 的核；设 \(p: E \to E/N\) 是典范满射，并设 \(v: E/N \to F\) 是 u 在商空间上诱导的映射。集 \(p(B)\) 是 E/N 中范数 \textless{} 1 的元素所成的集。我们有

\[
\mathrm{P} = u (\mathrm{E}) - u (\mathrm{B}) = v (\mathrm{E} / \mathrm{N}) - v (p (\mathrm{B})) = v ((\mathrm{E} / \mathrm{N}) - p (\mathrm{B}))
\]

因为映射 \(v\) 是单射。换言之，P 由 F 中形如 \(v(y)\) （其中 \(y \in \mathrm{E/N}\) 且 \(\|y\| \geqslant 1\) ）的元素组成。用 \(d_{\mathrm{P}}\) 表示在 F 中从 0 到 P 的距离。我们有

\[
d_{\mathrm{P}} = \inf_{\substack{y\in \mathrm{E} / \mathrm{N}\\ \| y\| \geqslant 1}}\| v(y)\| = \inf_{\substack{y\in \mathrm{E} / \mathrm{N}\\ \| y\| = 1}}\| v(y)\| = ((u))
\]

由 (6) 即得。

设 \(u\) 是严格态射。设 \(\varepsilon > 0\) 。集 \(\varepsilon u(\mathrm{B})\) 是 \(u(\mathrm{E})\) 中 0 的一个邻域。\(u(\mathrm{B})\) 在 \(u(\mathrm{E})\) 中的闭包等于 \(\overline{u(\mathrm{B})} \cap u(\mathrm{E})\) 。它包含于集 \(u(\mathrm{B}) + \varepsilon u(\mathrm{B})\) ，后者等于 \((1 + \varepsilon)u(\mathrm{B})\) ，因为 \(u(\mathrm{B})\) 是凸的。于是有 \((1 + \varepsilon)\mathrm{P} \subset \mathrm{Q} \subset \mathrm{P}\) ，且在 F 中从 0 到 Q 的距离 \(d_{\mathrm{Q}}\) 满足不等式 \(d_{\mathrm{P}} \leqslant d_{\mathrm{Q}} \leqslant (1 + \varepsilon)d_{\mathrm{P}}\) 。由于这对每个 \(\varepsilon > 0\) 都成立，故有 \(d_{\mathrm{Q}} = d_{\mathrm{P}} = ((u))\) 。

设 \(u\) 不是严格态射，但赋范空间 E 完备。此时有 \(((u)) = 0\) 。\(u(\mathrm{B})\) 在 \(u(\mathrm{E})\) 中的闭包（它等于 \(\overline{u(\mathrm{B})} \cap u(\mathrm{E})\) ）不是 \(u(\mathrm{E})\) 中 0 的邻域（EVT, I, p. 17, 定理 1）。因此存在 Q 中任意接近 0 的点，从而 \(d_{\mathrm{Q}} = 0 = ((u))\) 。

命题 7. --- 设 E 是巴拿赫空间，F 是赋范空间。映射 \(u \mapsto ((u))\) （从 \(\mathcal{L}(\mathrm{E};\mathrm{F})\) 到 \(\overline{\mathbf{R}}\) 中）是上半连续的。

设 \(u \in \mathcal{L}(\mathrm{E};\mathrm{F})\) 。只需证明：对每个实数 \(c > ((u))\) ，由元素 \(v \in \mathcal{L}(\mathrm{E};\mathrm{F})\) （满足 \(((v)) < c\) ）所成的集是 \(u\) 的一个邻域。用 B 表示 E 中范数 \(< 1\) 的元素所成的集。由命题 6，存在 \(y \in \mathrm{E}\) ，使得 \(u(y) \notin \overline{u(\mathrm{B})}\) 且 \(\| u(y)\| < c\) 。距离 \(d\) （从 \(u(y)\) 到闭集 \(\overline{u(\mathrm{B})}\) ）是严格正的。由元素 \(v\) （\(\mathcal{L}(\mathrm{E};\mathrm{F})\) 中的、满足关系式 \(\| v(y)\| < c\) 与 \(\| u - v\|(1 + \| y\|) < d\) 者）所成的集 V 是 \(u\) 在 \(\mathcal{L}(\mathrm{E};\mathrm{F})\) 中的一个邻域。设 \(v \in \mathrm{V}\) 。对每个 \(x \in \mathrm{B}\) ，我们有

\[
\begin{array}{c} d \leqslant \| u (y) - u (x) \| \leqslant \| v (y) - v (x) \| + \| u - v \| (\| y \| + \| x \|) \\ <   \| v (y) - v (x) \| + d. \end{array}
\]

于是 \(v(y)\) 不属于 \(v(\mathrm{B})\) ，且由命题 6 有 \(((v)) \leqslant \| v(y) \| < c\) 。

命题 8. --- 设 E 是巴拿赫空间，F 是赋范空间。对每个 \(u \in \mathcal{L}(\mathrm{E};\mathrm{F})\) ，我们有 \(((u)) = ((^t u))\) 。

用 \(j\) 表示赋范空间 F 到其完备化 \(\widehat{\mathrm{F}}\) 中的典范单射。我们有 \(((u)) = ((j \circ u))\) 。由于线性映射 \(^t j: (\widehat{\mathrm{F}})' \to \mathrm{F}'\) 是双射且等距，也有 \(((^t u)) = ((^t (j \circ u)))\) 。因此只需就赋范空间 F 完备的情形证明本命题，我们在证明的其余部分即作此假设。

若 \(u\) 为零，则有 \(((u)) = ((^t u)) = +\infty\) 。若 \(u\) 不是严格态射，则 \(^t u\) 也不是（EVT, IV, p. 29, 推论 3），且有 \(((u)) = ((^t u)) = 0\) 。

现在设 u 是非零的严格态射。用 N 表示 u 的核，I 表示它的像，并考察 u 的典范分解：

\[
\mathrm{E} \xrightarrow {p} \mathrm{E/N} \xrightarrow {v} \mathrm{I} \xrightarrow {i} \mathrm{F}.
\]

\(^{t}u\) 的核是 I 在 F' 中的正交 I°（EVT, IV, p. 27, 命题 2），且 \(^{t}i\) 通过过渡到商空间诱导一个从 F'/I° 到 I' 上的等距 \(\iota\) （EVT, IV, p. 9, 命题 10）。此外，\(^{t}p\) 给出从 (E/N)' 到 N 在 E' 中的正交 N° 上的一个等距（EVT, IV, p. 8, 命题 9）。因此 \(^{t}u\) 的典范分解是

\[
\mathrm{F} ^ {\prime} \longrightarrow \mathrm{F} ^ {\prime} / \mathrm{I} ^ {\circ} \stackrel {v ^ {\prime}} {\longrightarrow} \mathrm{N} ^ {\circ} \longrightarrow \mathrm{E} ^ {\prime},
\]

其中 \(v^{\prime} = \pi \circ{}^{t}v\circ \iota\)。此时有

\[
\| (v ^ {\prime}) ^ {- 1} \| = \| (^ {t} v) ^ {- 1} \| = \| ^ {t} (v ^ {- 1}) \| = \| v ^ {- 1} \|
\]

（EVT, IV, p. 7, 命题 8），从而由公式 (10) 得 \(((u)) = ((^t u))\) 。

\subsection{赋范空间的有限维向量子空间}

下面的命题将在 TA（待出版）中作为博苏克–乌拉姆定理的推论而证明。

定理 3. — 设 n 是正整数，V 是维数为 \(n + 1\) 的实赋范向量空间，W 是 V 的维数为 n 的向量子空间。设 S 是 V 的单位球面。不存在连续映射 \(f: S \to W\) ，使得 \(\|f(x) - x\| < 1\) 对一切 \(x \in S\) 成立。

定理 4（克赖因、克拉索塞利斯基、米尔曼）. — 设 E 是赋范空间，F 和 G 是 E 的向量子空间。假设 G 的维数有限且严格小于 F 的维数。则存在 F 中范数为 1 的元素，它到 G 的距离等于 1。

只需处理域 K 等于 R 的情形。设 n 是 G 的维数。把 F 换成 F 的一个包含 G 且维数为 \(n+1\) 的向量子空间，就化归于 \(\dim(F)=n+1\) 的情形。用反证法，假设定理的结论不成立。设 S 是 F 的单位球面。对每个 \(x\in S\) ，x 到 G 的距离严格小于 \(\|x\|=1\) ，故可选取 G 的元素 \(y(x)\) ，使得 \(\|x-y(x)\|<1\) 。用 \(V_{x}\) 表示 S 中满足 \(\|z-y(x)\|<1\) 的元素 z 所成的集；它是 x 在 S 中的一个开邻域。存在 S 上的局部有限连续单位分解 \((\varphi_{x})_{x\in\mathrm{S}}\) ，它从属于 S 的覆盖 \((\mathrm{V}_{x})_{x\in\mathrm{S}}\) （TG, IX, p. 46, 命题 3 及 p. 51, 命题 6）。用 \(f\colon S\to G\) 表示由下式给出的连续映射

\[
f (z) = \sum_ {x \in \mathrm{S}} \varphi_ {x} (z) y (x).
\]

对一切 \(z \in S\)，我们有 \(\varphi_x(z) \geqslant 0\) （对一切 \(x \in S\)），且存在 \(x \in S\) 使 \(\varphi_x(z) > 0\)，因为 \(\sum_{x \in S} \varphi_x(z) = 1\)，由此

\[
\begin{array}{c} \| z - f (z) \| = \left\| \sum_ {x \in \mathrm{S}} \varphi_ {x} (z) (z - y (x)) \right\| \leqslant \sum_ {x \in \mathrm{S}} \varphi_ {x} (z) \| z - y (x) \| \\ <   \sum_ {x \in \mathrm{S}} \varphi_ {x} (z) = 1. \end{array}
\]

f 的这一性质与定理 3 矛盾。

命题 9. — 设 E 和 F 是赋范空间，\(n \in N\)，并设 u 是从 E 到 F 中的连续线性映射，其核的维数为 n。u 的余范数等于从 u 到映射 \(v \in \mathcal{L}(\mathrm{E};\mathrm{F})\) （其核的维数至少为 \(n + 1\)）所成的集的距离 d。

当 \(u = 0\) 时，我们有 \(((u)) = +\infty\) 且 \(\dim(\mathrm{E}) = n\) ，从而 \(d = +\infty\) 。以下设 \(u\) 非零，于是 \(((u)) < +\infty\) 。设 \(v\) 是 \(\mathcal{L}(\mathrm{E};\mathrm{F})\) 中满足 \(\| u - v \| < ((u))\) 的元素，我们来证明它的核的维数 \(\leqslant n\) 。设 \(x\) 是 \(\operatorname{Ker}(v)\) 的一个范数为 1 的元素，\(y\) 是它在 \(\operatorname{E}/\operatorname{Ker}(u)\) 中的像。我们有（III, p. 61 的公式 (5)）

\[
((u)) \| y \| \leqslant \| u (x) \| = \| (u - v) (x) \| \leqslant \| u - v \| <   ((u))
\]

从而 \(\|y\| < 1\) 。而 \(\|y\|\) 正是 x 到 \(\operatorname{Ker}(u)\) 的距离。于是定理 4 蕴含 \(\dim \operatorname{Ker}(v) \leqslant \dim \operatorname{Ker}(u) = n\) 。这就证明了不等式 \(((u)) \leqslant d\) 。反向不等式 \(d \leqslant ((u))\) 由下面更精确的引理得出。

引理 3. — 设 c 是满足 \(c > ((u))\) 的实数。存在连续线性映射 \(h: E \to F\) ，其秩为 1、范数 \textless{} c，使得 \(u + h\) 的核包含 u 的核且维数为 \(n + 1\) 。

用 \(p\) 表示从 E 到 \(\mathrm{E} / \mathrm{Ker}(u)\) 上的典范映射。存在 \(a \in \mathrm{E}\) ，使得 \(\| p(a) \| = 1\) 且 \(\| u(a) \| < c\) （III, p. 61 的公式 (6)）。由哈恩–巴拿赫定理（EVT, II, p. 24, 推论 2），存在连续线性泛函 \(f\) （在赋范空间 \(\mathrm{E} / \mathrm{Ker}(u)\) 上），使得 \(f(p(a)) = 1\) 且 \(\| f \| = 1\) 。线性映射 \(h: x \mapsto -(f \circ p)(x)u(a)\) （从 E 到 F 中）连续且秩为 1，并且有 \(\| h \| \leqslant \| f \| \| p \| \| u(a) \| < c\) 。映射 \(u + h\) 的核包含 \(u\) 的核和 \(a\) ；由于 \(a \notin \mathrm{Ker}(u)\) ，它的维数至少为 \(n + 1\) 。另一方面，线性映射 \(u\) 通过过渡到子空间诱导一个从 \(\mathrm{Ker}(u + h)\) 到 \(\mathrm{Im}(h)\) 中的线性映射，其核为 \(\mathrm{Ker}(u) \cap \mathrm{Ker}(u + h)\) ，于是 \(\dim(\mathrm{Ker}(u + h)) \leqslant \dim(\mathrm{Ker}(u)) + 1\) 。结论得证。

推论 1. --- 设 E 和 F 是赋范空间，u、v 是从 E 到 F 中的非零连续线性映射，其核有相同的有限维数。则我们有

\[
\left| \left((u)\right) - ((v)) \right| \leqslant \| u - v \|.
\]

用 \(n\) 表示 \(u\) 与 \(v\) 的核的共同维数。用 A 表示从 E 到 F 的、核的维数 \(\geqslant n + 1\) 的连续线性映射所成的集。由于 \(u \neq 0\) ，我们有 \(\dim(\mathrm{E}) > n\) ，且集 A 包含 0。由命题 9，\(((u))\) 与 \(((v))\) 分别是 \(u\) 与 \(v\) 到 \(\mathcal{L}(\mathrm{E};\mathrm{F})\) 中集 A 的距离。于是只需应用公式 \(|d(u,\mathrm{A}) - d(v,\mathrm{A})| \leqslant \|u - v\|\) （参见 TG, IX, p. 13）。

推论 2. --- 设 E 和 F 是巴拿赫空间，u 和 v 是从 E 到 F 中的非零连续线性映射，其像在 F 中有相同的有限余维数。则有

\[
| ((u)) - ((v)) | \leqslant \| u - v \|.
\]

态射 u 是严格的（III, p. 52, 引理 6）。连续线性映射 \({}^{t}u\) 的核是正交 \((\mathrm{Im}(u))^{\circ}\) （\(\mathrm{Im}(u)\) 的正交）（EVT, IV, p. 27, 命题 2），由此 \(\dim(\mathrm{Ker}({}^{t}u)) = \mathrm{codim}_{\mathrm{F}}(\mathrm{Im}(u))\) ，类似地 \(\dim(\mathrm{Ker}({}^{t}v)) = \mathrm{codim}_{\mathrm{F}}(\mathrm{Im}(v))\) 。因此 \({}^{t}u\) 与 \({}^{t}v\) 的核有相同的有限维数。由于

\[
\| ^ {t} u - ^ {t} v \| = \| u - v \|, \quad ((^ {t} u)) = ((u)), \quad ((^ {t} v)) = ((v))
\]

（EVT, IV, p. 7, 命题 8 及 III, p. 63, 命题 8），把推论 1 应用于 \(^{t}u\) 和 \(^{t}v\) 即得该断言。

\subsection{连续单射或满射线性映射的扰动}

在本节中，我们采用下列约定：若 E 是有限维向量空间，\(\dim(\mathrm{E})\) 表示它的维数；若 E 是无穷维向量空间，则令 \(\dim(\mathrm{E}) = +\infty \in \overline{\mathbf{R}}\) 。若 u 是核或余核为有限维的线性映射，则令 \(\operatorname{ind}(u) = \dim \operatorname{Coker}(u) - \dim \operatorname{Ker}(u)\) ，其中的运算在 \(\overline{R}\) 中进行。

设 E 和 F 是赋范空间。用 \(\mathcal{M}(\mathrm{E};\mathrm{F})\) 表示从 E 到 F 中的严格单形态射所成的集，用 \(\mathcal{Q}\mathcal{M}(\mathrm{E};\mathrm{F})\) 表示从 E 到 F 中的、核为有限维的严格态射所成的集。

命题 10. --- 设 E 和 F 是赋范空间。集 \(\mathcal{M}(\mathrm{E};\mathrm{F})\) 在 \(\mathcal{L}(\mathrm{E};\mathrm{F})\) 中是开的。它是 \(\mathcal{L}(\mathrm{E};\mathrm{F})\) 中单射映射所成的集的内部。

用 A 表示从 E 到 F 中的连续单射线性映射所成的集。为使映射 \(u \in A\) 是严格态射，必须且只需它的余范数 \(((u))\) 严格为正（III, p. 62, 注记）。而 \(((u))\) 是从 \(u\) 到 A 在 \(\mathcal{L}(\mathrm{E};\mathrm{F})\) 中的补集的距离（III, p. 65, 命题 9）。命题得证。

命题 11. --- 设 E 和 F 是巴拿赫空间。集 \(\mathcal{Q}\mathcal{M}(\mathrm{E};\mathrm{F})\) 在 \(\mathcal{L}(\mathrm{E};\mathrm{F})\) 中是开的。它是 \(\mathcal{L}(\mathrm{E};\mathrm{F})\) 中核为有限维的映射所成的集的内部。

用 A 表示从 E 到 F 中的、核为有限维的连续线性映射所成的集。

设 u 是 \(\mathcal{Q}\mathcal{M}(\mathrm{E};\mathrm{F})\) 的一个元素。此时有 \(((u)) > 0\) （III, p. 62, 注记）。每个元素 \(v \in \mathcal{L}(\mathrm{E};\mathrm{F})\) （与 u 的距离 \(<((u))\) ）都属于 A（III, p. 65, 命题 9），故 u 是 A 的一个内点。

反之，设 \(u\) 是 A 的一个不是严格态射的元素。我们有 \(((u)) = 0\) （III, p. 62, 注记）。设 \(v\) 是 \(\mathcal{L}(\mathrm{E};\mathrm{F})\) 的一个与 \(u\) 相差一个有限秩线性映射的元素；由 III, p. 40 的推论 1，态射 \(v\) 不可能是具有闭像的严格态射；由于从 E 到 F 中的严格态射在 F 中有闭像（EVT, IV, p. 28, 定理 1），这意味着 v 不是严格态射。于是有 \(((v)) = 0\) 。设 \(\varepsilon\) 为一实数，\textgreater{} 0。借助前面的考察，III, p. 66 的引理 3 使我们可以归纳地构造元素的一个序列 \((u_{m})_{m \in \mathbf{N}}\) （\(\mathcal{L}(\mathrm{E}; \mathrm{F})\) 中的），它满足下列条件：

\begin{enumerate}
\def\labelenumi{(\roman{enumi})}
\item
  有 \(u_{0}=u\) ；
\item
  对每个 \(m \geqslant 0\) ，\(u_{m+1} - u_m\) 是秩为 1 且范数 \(\leqslant 2^{-m-1}\varepsilon\) 的连续线性映射；
\item
  对每个 \(m \geqslant 0\) ，\(u_m\) 的核维数为 \(n + m\) ，且包含于 \(u_{m+1}\) 的核中。
\end{enumerate}

序列 \((u_{m})\) 是巴拿赫空间 \(\mathcal{L}(\mathrm{E};\mathrm{F})\) 中的柯西序列。设 \(u'\) 是它的极限。\(u'\) 的核包含 \(u_{m}\) 的核（对一切 \(m \geqslant 0\) ）；它是无穷维的。由于

\[
\left\| u ^ {\prime} - u \right\| \leqslant \sum_ {m = 0} ^ {\infty} \left\| u _ {m + 1} - u _ {m} \right\| \leqslant \varepsilon \sum_ {m = 0} ^ {\infty} 2 ^ {- m - 1} = \varepsilon ,
\]

u 到 A 的补集的距离小于 \(\varepsilon\) 。由于这对每个 \(\varepsilon > 0\) 都成立，由此得出 u 不是 A 的内点。

命题 12. --- 设 E 和 F 是巴拿赫空间，u 是 \(\mathcal{Q}\mathcal{M}(\mathrm{E};\mathrm{F})\) 的一个元素。每个 \(v \in \mathcal{L}(\mathrm{E};\mathrm{F})\) （满足 \(\|v - u\| < ((u))\) ）都属于 \(\mathcal{Q}\mathcal{M}(\mathrm{E};\mathrm{F})\) 且满足关系式

\[
\begin{array}{c} \dim \operatorname{Ker} (v) \leqslant \dim \operatorname{Ker} (u), \\ \dim \operatorname{Coker} (v) \leqslant \dim \operatorname{Coker} (u), \\ \operatorname{ind} (v) = \operatorname{ind} (u). \end{array}
\]

由于 \(u\) 是严格的，我们有 \(((u)) > 0\) 。用 B 表示以 \(v\) 为中心、\(((u))\) 为半径的开球（在 \(\mathcal{L}(\mathrm{E};\mathrm{F})\) 中）。对每个 \(v \in \mathrm{B}\) ，我们有 \(\dim \operatorname{Ker}(v) \leqslant \dim \operatorname{Ker}(u)\) （III, p. 65, 命题 9）且 \(v \in \mathcal{Q}\mathcal{M}(\mathrm{E};\mathrm{F})\) （命题 11）。

对 \(r \in \mathbf{Z} \cup \{+\infty\}\) ，用 \(B_r\) 表示由元素 \(v \in B\) （满足 \(\text{ind}(v) = r\) ）所成的集。若 \(r \in \mathbf{Z}\) ，则集 \(B_r\) 是从 E 到 F 的、指标为 \(r\) 且属于 B 的弗雷德霍姆映射所成的集（III, p. 52, 命题 11）；由 III, p. 58 的定理 1，它在 B 中是开的。

我们来证明诸集 \(B_{r}\) 在 B 中是闭的。设 \(v \in B\) 是附着于 \(B_{r}\) 的一点。由于诸集 \(B_{s}\) （\(s \in Z\) ）在 B 中是开的且两两不相交，我们有 \(v \in B_{r}\) 或 \(v \in B_{+\infty}\) 。

设 \(v \in B_{+\infty}\) 。用 n 表示 \(\operatorname{Ker}(u)\) 的维数。选取 F 的一个维数为 \(r + 2n + 1\) 的向量子空间 T，它与 \(\operatorname{Im}(v)\) 的交缩减为 0；再选取 \(\operatorname{Ker}(v)\) 在 E 中的一个拓扑补 S（III, p. 55, 命题 1）。考察连续线性映射 \(f: (s, t) \mapsto v(s) + t\) （从 S × T 到 F）。它是单射的。线性映射 v 是严格态射，因为 v 属于 B ⊂ \(\mathcal{Q}\mathcal{M}(\mathrm{E};\mathrm{F})\) 。于是 v 的像是闭的（EVT, IV, p. 28, 定理 1）。由 III, p. 39 的命题 1，v 在 S 上的限制是从 S 到 F 的具有闭像的严格态射，而 f 是从 S × T 到 F 的具有闭像的严格态射。

由命题 10，存在 \(v\) 在 B 中的一个邻域 U，使得对每个 \(w \in \mathrm{U}\) ，线性映射 \((s,t) \mapsto w(s) + t\) （从 \(\mathrm{S} \times \mathrm{T}\) 到 F 中）是单射的。设 \(w\) 是 U 的一个元素。此时有 \(\operatorname{Ker}(w) \cap \mathrm{S} = \{0\}\) ，从而 \(w\) 通过限制并过渡到商空间诱导一个从 \(\operatorname{Ker}(w)\) 到 E/S 中的单射线性映射，这蕴含 \(\operatorname{Ker}(w)\) 的维数至多为 \(n\) 。我们还有 \(w(\mathrm{S}) \cap \mathrm{T} = \{0\}\) ，由此

\[
\operatorname{codim} _ {\mathrm{F}} (\operatorname{Im} (w)) \geqslant \dim (\mathrm{T}) - \operatorname{codim} _ {\mathrm{E}} (\mathrm{S}) = r + n + 1,
\]

这蕴含 \(\operatorname{ind}(w) \geqslant r + 1\) ，且与 v 附着于 B \(_{r}\) 这一假设相矛盾。

若 \(r\) 是 \(\mathbf{Z}\) 中使 \(B_r\) 非空的元素，则我们有 \(B_r = B\) ，因为 \(B_r\) 在 \(B\) 中既开又闭且 \(B\) 连通。若 \(B_r\) 对每个 \(r \in \mathbf{Z}\) 都是空的，则有 \(\text{ind}(v) = +\infty\) （对一切 \(v \in B\) ）。于是映射 \(v \mapsto \text{ind}(v)\) 在 \(B\) 上是常数。最后，对每个 \(v \in B\) ，我们有

\[
\begin{array}{c} \dim \operatorname{Coker} (v) = \operatorname{ind} (v) + \dim \operatorname{Ker} (v) \\ \leqslant \operatorname{ind} (u) + \dim \operatorname{Ker} (u) = \dim \operatorname{Coker} (u). \end{array}
\]

证明到此结束。

推论 1. --- 由 \(u \mapsto \dim \operatorname{Ker}(u)\) 与 \(u \mapsto \dim \operatorname{Coker}(u)\) 在 \(\mathcal{Q}\mathcal{M}(\mathrm{E};\mathrm{F})\) 上定义的函数是上半连续的。函数 \(u \mapsto \operatorname{ind}(u)\) 在 \(\mathcal{Q}\mathcal{M}(\mathrm{E};\mathrm{F})\) 上是局部常数的。

推论 2. --- 函数 \(u \mapsto \dim \operatorname{Coker}(u)\) 在从 E 到 F 中的严格单形态射所成的集 \(\mathcal{M}(\mathrm{E};\mathrm{F})\) 上是局部常数的。

事实上，有 \(\dim \operatorname{Coker}(u) = \operatorname{ind}(u)\) （对 \(u \in \mathcal{M}(\mathrm{E};\mathrm{F})\) ）。

引理 4. --- 设 E 和 F 是巴拿赫空间。为使 \(\mathcal{L}(\mathrm{E};\mathrm{F})\) 的元素 u 是从 E 到 F 中的严格态射，必须且只需 \(^{t}u\) 是从 F' 到 E' 中的严格态射。此时有 \(\dim \operatorname{Coker}(^{t}u)=\dim\operatorname{Ker}(u)\) 与 \(\dim\operatorname{Ker}(^{t}u)=\dim\operatorname{Coker}(u)\) 。

为使 u 是严格态射，必须且只需 \(^{t}u\) 是（EVT, IV, p. 29, 推论 3）；此时 u 的像是闭的（EVT, IV, p. 28, 定理 1），且向量空间 \(\mathrm{Ker}(^{t}u)\) （相应地，\(\mathrm{Coker}(^{t}u)\) ）典范地同构于赋范空间 \(\mathrm{Coker}(u)\) （相应地，\(\mathrm{Ker}(u)\) ）的对偶（依照 EVT, IV, p. 27, 命题 2）。引理得证。

设 E 和 F 是巴拿赫空间。用 \(\mathcal{E}(\mathrm{E};\mathrm{F})\) 表示从 E 到 F 中的连续满射线性映射所成的集，用 \(\mathcal{Q}\mathcal{E}(\mathrm{E};\mathrm{F})\) 表示从 E 到 F 中的、像有有限余维数的连续线性映射所成的集。\(\mathcal{Q}\mathcal{E}(\mathrm{E};\mathrm{F})\) 的每个元素都是具有闭像的严格态射（III, p. 52, 引理 6）。由引理 4 得出，\(\mathcal{E}(\mathrm{E};\mathrm{F})\) 与 \(\mathcal{Q}\mathcal{E}(\mathrm{E};\mathrm{F})\) 分别是 \(\mathcal{M}(\mathrm{F}^{\prime};\mathrm{E}^{\prime})\) 与 \(\mathcal{Q}\mathcal{M}(\mathrm{F}^{\prime};\mathrm{E}^{\prime})\) 在连续映射 \(u\mapsto{}^{t}u\) （从 \(\mathcal{L}(\mathrm{E};\mathrm{F})\) 到 \(\mathcal{L}(\mathrm{F}^{\prime};\mathrm{E}^{\prime})\) 中）下的原像。

命题 13. --- 设 E 和 F 是巴拿赫空间。集 \(\mathcal{E}(\mathrm{E};\mathrm{F})\) 与 \(\mathcal{Q}\mathcal{E}(\mathrm{E};\mathrm{F})\) 在 \(\mathcal{L}(\mathrm{E};\mathrm{F})\) 中是开的。更精确地说，若 \(u\) 是 \(\mathcal{Q}\mathcal{E}(\mathrm{E};\mathrm{F})\) 的元素，而 \(v\) 是 \(\mathcal{L}(\mathrm{E};\mathrm{F})\) 中满足 \(\| v - u\| < ((u))\) 的元素，则我们有 \(v \in \mathcal{Q}\mathcal{E}(\mathrm{E};\mathrm{F})\) ，且

\(\dim \operatorname{Ker}(v) \leqslant \dim \operatorname{Ker}(u), \quad \dim \operatorname{Coker}(v) \leqslant \dim \operatorname{Coker}(u),\)

\[
\operatorname{ind} (v) = \operatorname{ind} (u).
\]

我们前面已经看到 \(^{t}u \in \mathcal{Q}\mathcal{M}(F';E')\) 。此外，对每个元素 \(v\) （在 \(\mathcal{L}(E;F)\) 中），有 \(\|^{t}v - ^{t}u\| = \|v - u\|\) 与 \(((^{t}u)) = ((u))\) （EVT, IV, p. 7, 命题 8 及 III, p. 63, 命题 8）。由命题 12，从这些关系式得出：若 \(\|v - u\| < ((u))\) ，则 \(^{t}v\) 属于 \(\mathcal{Q}\mathcal{M}(F';E')\) ，且有不等式

\[
\dim \operatorname{Ker} (^ {t} v) \leqslant \dim \operatorname{Ker} (^ {t} u), \quad \dim \operatorname{Coker} (^ {t} v) \leqslant \dim \operatorname{Coker} (^ {t} u)
\]

以及等式 \(\mathrm{ind}(^{t}v)=\mathrm{ind}(^{t}u)\) 。于是命题由引理 4 得出。

推论 1. --- 由 \(u \mapsto \dim \operatorname{Ker}(u)\) 与 \(u \mapsto \dim \operatorname{Coker}(u)\) 在 \(\mathcal{Q}\mathcal{E}(\mathrm{E};\mathrm{F})\) 上定义的函数是上半连续的。函数 \(u \mapsto \operatorname{ind}(u)\) 在 \(\mathcal{Q}\mathcal{E}(\mathrm{E};\mathrm{F})\) 上是局部常数的。

推论 2. --- 由 \(u \mapsto \dim \operatorname{Ker}(u)\) 定义的函数在 \(\mathcal{E}(\mathrm{E};\mathrm{F})\) 上是局部常数的。

事实上，有 \(\dim \operatorname{Ker}(u) = -\operatorname{ind}(u)\) （对一切 \(u \in \mathcal{E}(\mathrm{E};\mathrm{F})\) ）。

\section{借紧线性映射的扰动}

\subsection{严格态射与真性}

引理 1. --- 设 E 和 F 是拓扑向量空间，u 是从 E 到 F 中的连续线性映射，U 是 E 中 0 的一个邻域。假设 u 诱导一个从 U 到 F 的闭子集上的同胚。则 u 的像 I 是闭的，且 u 诱导一个从 E 到 I 上的同胚。

集 \(u(\mathrm{U})\) 是 I 中 0 的一个邻域；它在 F 中是闭的，故 F 的子群 I 在 0 处局部闭，从而它是闭的（TG, III, p. 7, 命题 4）。

由于 \(\operatorname{Ker}(u)\) 与 U 只交于 0 这一点，有 \(\operatorname{Ker}(u) = \{0\}\) ，故映射 \(u\) 是单射。设 V 是包含于 U 的 E 中 0 的闭邻域。由于 \(u(\mathring{\mathrm{V}})\) 是 \(u(\mathrm{V})\) 中 0 的一个邻域，存在 F 中 0 的均衡邻域 W，使得 \(u(\mathrm{V}) \cap \mathrm{W} \subset u(\mathring{\mathrm{V}})\) 。集 \(u^{-1}(\mathrm{W})\) 是均衡的，故是连通的。它包含 0 且包含于 \(\mathring{\mathrm{V}} \cup (\mathrm{E} - \mathrm{V})\) 。由于 \(\mathring{\mathrm{V}}\) 与 E-V 是 E 的不相交开子集，集 \(u^{-1}(\mathrm{W})\) 包含于 \(\mathring{\mathrm{V}}\) ，从而 \(\mathrm{W} \cap \mathrm{I} \subset u(\mathring{\mathrm{V}})\) 。于是 \(u(\mathring{\mathrm{V}})\) 是 I 中 0 的一个邻域。这蕴含 \(u\) 诱导一个从 E 到 I 上的同胚。

命题 1. --- 设 E 是分离局部凸空间，F 是局部凸空间，u 是从 E 到 F 中的连续线性映射。下列条件等价：

\begin{enumerate}
\def\labelenumi{(\roman{enumi})}
\item
  映射 u 是严格态射，它的核有限维且它的像在 F 中是闭的；
\item
  存在 E 中 0 的闭邻域 V，使得 u 在 V 上的限制是真映射（TG, I, p. 72）。
\item
  \(\Longrightarrow\) (ii)：假设 u 满足条件 (i)。由于 u 的核有限维，它在 E 中有一个拓扑补 \(E_{1}\) （III, p. 55, 命题 1），且在 \(\operatorname{Ker}(u)\) 中存在 0 的一个紧邻域 C。把 E 与 \(E_{1} \times \operatorname{Ker}(u)\) 等同起来；此时集 \(V = E_{1} \times C\) 是 E 中 0 的一个闭邻域。u 在 V 上的限制是 \(E_{1} \times C\) 到 \(E_{1}\) 上的投影（它是真映射，TG, I, p. 77, 推论 5）与 \(u_{1}\) （即 u 在 \(E_{1}\) 上的限制）的复合。而 \(u_{1}\) 是从 \(E_{1}\) 到 F 的闭子空间上的一个同胚，故是真映射（TG, I, p. 72, 命题 2）。两个真映射的复合是真映射（TG, I, p. 73, 命题 5, a)），所以 u 在 V 上的限制是真映射。
\item
  \(\Longrightarrow\) (i)：设 V 是 E 中 0 的一个闭邻域，使得 u 在 V 上的限制 v 是真映射。此时集 \(\mathrm{V} \cap \mathrm{Ker}(u) = \overline{v}(\{0\})\) 是紧的（TG, I, p. 75, 定理 1）；于是向量空间 \(\mathrm{Ker}(u)\) 是局部紧的，从而有限维（EVT, I, p. 15, 定理 3）。设 \(E_{1}\) 是 \(\mathrm{Ker}(u)\) 在 E 中的一个拓扑补（III, p. 55 的命题 1）；令 \(V_{1} = E_{1} \cap V\) 。集 \(V_{1}\) 在 V 中是闭的。映射 \(u|V_{1}\) 是真映射（TG, I, p. 74, 推论 1）且是单射，故是从 \(V_{1}\) 到 F 的闭子集上的一个同胚（TG, I, p. 72, 命题 2）。由引理 1，u 在 \(E_{1}\) 上的限制是从 \(E_{1}\) 到 F 的闭子空间上的一个同胚，故 u 是具有闭像的严格态射。
\end{enumerate}

\subsection{单射或满射线性映射的扰动}

定理 1. --- 设 E 和 F 是分离局部凸空间，u 是从 E 到 F 中的严格态射，其核有限维且其像是闭的。设 h 是从 E 到 F 中的紧线性映射。线性映射 \(u + h\) 是严格态射，它的核有限维，且它的像是闭的。

由 III, p. 71 的命题 1，存在 E 中 0 的一个闭邻域 V，使得 \(u\) 在 V 上的限制是真映射。由于 \(h\) 是紧的，存在包含于 V 的 0 的闭邻域 W，使得集 \(h(\mathrm{W})\) 是相对紧的。令 \(\mathrm{C} = \overline{h(\mathrm{W})}\) 。\(u\) 在 W 上的限制是真映射（TG, I, p. 74, 推论 1）。映射 \(\alpha: x \mapsto (u(x), h(x))\) （从 W 到 \(\mathrm{F} \times \mathrm{C}\) 中）是真映射，因为映射 \(\mathrm{pr}_1 \circ \alpha = u|\mathrm{W}\) （从 W 到 F 中）是真映射（TG, I, p. 73, 命题 5, \(d\) ）。映射 \(\beta: (y, z) \mapsto (y + z, z)\) （从 \(\mathrm{F} \times \mathrm{C}\) 到 \(\mathrm{F} \times \mathrm{C}\) 中）是同胚，而映射 \(\mathrm{pr}_1\) （从 \(\mathrm{F} \times \mathrm{C}\) 到 F 中）是真映射（TG, I, p. 77, 推论 5）。于是复合映射 \(\mathrm{pr}_1 \circ \beta \circ \alpha\) （从 W 到 F 中）是真映射（TG, I, p. 73, 命题 5, \(a\) ）。但这个映射正是 \(u + h\) 在 W 上的限制。由 III, p. 71 的命题 1，\(u + h\) 是严格态射，它的核有限维，且它的像是闭的。

引理 2. --- 设 E 和 F 是弗雷歇空间，且 \(u \in \mathcal{L}(\mathrm{E};\mathrm{F})\) 。下列条件等价：

\begin{enumerate}
\def\labelenumi{(\roman{enumi})}
\item
  u 的余核有限维；
\item
  \(^{t}u \colon \mathrm{F}_c' \to \mathrm{E}_c'\) 是具有闭像且核有限维的严格态射。
\item
  \(\Longrightarrow\) (ii)：假设 u 的余核有限维。此时映射 u 是严格态射（III, p. 52 的引理 6）。由 EVT, IV, p. 27 的命题 2，\(^{t}u\) 的像在赋有弱拓扑的 E' 中是闭的，从而在 \(E_{c}'\) 中更是闭的。\(^{t}u\) 的核是 u 的像的正交补（同前）；故它有限维。最后，由 EVT, IV, p. 28 的定理 1，\(^{t}u\) 是从 \(F_{c}'\) 到 \(E_{c}'\) 中的严格态射。
\item
  \(\Longrightarrow\) (i)：假设 \(^t u\colon F_c' \to E_c'\) 是严格态射。由 EVT, IV, p. 28 的定理 1，\(u\) 的像是闭的。\(^t u\) 的核是 \(\operatorname{Im}(u)\) 的正交补（依照 EVT, IV, p. 27, 命题 2）；若这个核有限维，则 \(u\) 的像在 F 中有有限余维数。
\end{enumerate}

定理 2. --- 设 E 和 F 是弗雷歇空间，u: E → F 是余核有限维的连续线性映射，h: E → F 是紧线性映射。线性映射 u + h 是严格态射，它的余核有限维，且它的像是闭的。

由 III, p. 52 的引理 6，只需证明 \(u + h\) 的余核有限维。而由 III, p. 9 的命题 9，映射 \(^{t}h\) （从 \(F_{c}^{\prime}\) 到 \(E_{c}^{\prime}\) 中）是紧的。定理 2 由定理 1 与引理 2 得出。

\subsection{弗雷德霍姆映射的扰动}

定理 3. --- 设 E 是局部凸空间，F 是分离局部凸空间，u 是从 E 到 F 中的弗雷德霍姆映射，h 是从 E 到 F 中的紧线性映射。则 \(u + h\) 是弗雷德霍姆算子，且有 \(\operatorname{ind}(u + h) = \operatorname{ind}(u)\) 。

我们将分三步证明该定理。

\begin{enumerate}
\def\labelenumi{\Alph{enumi})}
\item
  假设 E 和 F 是巴拿赫空间。III, p. 42 的命题 2 表明线性映射 u 是具有闭像且核与余核均有限维的严格态射。由于 E 和 F 是巴拿赫空间，III, p. 73 的定理 1 与定理 2 蕴含：对每个 \(t \in [0,1]\) ，线性映射 \(u_t = u + th\) 具有这些同样的性质，从而是弗雷德霍姆算子（III, p. 42, 命题 2）。映射 \(t \mapsto u_t\) （从 \([0,1]\) 到从 E 到 F 的弗雷德霍姆映射所成的集 \(\mathcal{F}(\mathrm{E};\mathrm{F})\) 中）是连续的。由 III, p. 58 的定理 1，映射 \(t \mapsto \operatorname{ind}(u_t)\) 是局部常数的。由于区间 \([0,1]\) 连通，这个映射是常数。于是我们有 \(\operatorname{ind}(u) = \operatorname{ind}(u + h)\) 。
\item
  我们假设 E = F，\(u = 1_{E}\) 。由 III, p. 5 的命题 5，存在一个巴拿赫空间 G、一个连续线性映射 \(p: E \to G\) 和一个紧线性映射 \(q: G \to E\)，使得 \(h = q \circ p\) 。G 的自同态 \(p \circ q\) 是紧的。由 A)，\(1_{G} + p \circ q\) 是 G 的指标为零的弗雷德霍姆自同态。但此时 \(1_{E} + h = 1_{E} + q \circ p\) 是 E 的指标为零的弗雷德霍姆自同态（III, p. 49, 命题 10, f)）。
\item
  一般情形。设 \(v\) 是 \(u\) 的一个拟逆。自同态 \(u \circ v\) 与 \((u + h) \circ v\) （\(F\) 的）与 \(1_{F}\) 都只差一个紧线性映射。由 B)，它们是 \(F\) 的指标为零的弗雷德霍姆自同态。由此得出 \(u + h\) 是弗雷德霍姆自同态（III, p. 40, 第 2 目）且与 \(u\) 有相同的指标（III, p. 43, 第 3 目, 公式 (2)）。
\end{enumerate}

推论 1. --- 设 E 和 F 是分离局部凸空间，且 \(u \in \mathcal{L}(\mathrm{E};\mathrm{F})\) 。为使 u 是弗雷德霍姆算子，必须且只需存在一个映射 \(v \in \mathcal{L}(\mathrm{F};\mathrm{E})\) ，使得线性映射 \(1_{\mathrm{E}} - v \circ u\) 与 \(1_{\mathrm{F}} - u \circ v\) 都是紧的。

由于有限秩的连续线性映射是紧的，所述条件是必要的。

设 \(v\) 是 \(\mathcal{L}(\mathrm{F};\mathrm{E})\) 中使得线性映射 \(1_{\mathrm{E}} - v \circ u\) 与 \(1_{\mathrm{F}} - u \circ v\) 为紧的的元素。由定理 3，\(v \circ u\) 与 \(u \circ v\) 是弗雷德霍姆映射。设 \(p\) 与 \(q\) 分别是 \(v \circ u\) 与 \(u \circ v\) 的拟逆。令 \(w = p \circ v\) ，\(w' = v \circ q\) 。采用 III, p. 40 第 2 目的记号，我们有 \(w \circ u \equiv 1_{\mathrm{E}}\) 与 \(u \circ w' \equiv 1_{\mathrm{F}}\) ，由此 \(w \equiv w \circ u \circ w' \equiv w'\) 。于是得出 \(w\) 是 \(u\) 的一个拟逆。

推论 2. --- 设 E 和 F 是分离局部凸空间，且 \(u \in \mathcal{L}(\mathrm{E};\mathrm{F})\) 。下列条件等价：

\begin{enumerate}
\def\labelenumi{(\roman{enumi})}
\item
  映射 u 是指标为零的弗雷德霍姆映射；
\item
  存在一个从 E 到 F 上的同构 v，使得 u - v 有限秩；
\item
  存在同构 \(v\) （从 \(E\) 到 \(F\) 上的），使得 \(u - v\) 是紧的。
\item
  \(\Longrightarrow\) (ii)：假设 u 是指标为零的弗雷德霍姆算子。存在拓扑直和分解 \(E = E_{1} \oplus E_{2}\) 与 \(F = F_{1} \oplus F_{2}\) （其中 \(E_{2}\) 与 \(F_{2}\) 有限维），使得 u 在 \(E_{2}\) 上为零，且通过限制定义一个同构 \(u_{1}\) （从 \(E_{1}\) 到 \(F_{1}\) 上）（III, p. 42, 命题 2）。若 \(\text{ind}(u) = 0\) ，则 \(E_{2}\) 与 \(F_{2}\) 的维数相等，且存在从 E 到 F 上的一个同构 v，它是 \(u_{1}\) 的扩张。u - v 的核包含 \(E_{1}\) ，故 u - v 有限秩。
\end{enumerate}

由于从 E 到 F 中的每个有限秩连续线性映射都是紧的，我们有 (ii) \(\Longrightarrow\) (iii)；而由定理 3 有 (iii) \(\Longrightarrow\) (i)。

\subsection{里斯自同态的扰动}

设 E 是巴拿赫空间。回忆（参见 III, p. 5, 命题 3）：E 的紧自同态所成的集 \(\mathcal{L}^c (\mathrm{E})\) 构成巴拿赫代数 \(\mathcal{L}(\mathrm{E})\) 的一个闭双边理想。商巴拿赫代数称为 E 的卡尔金代数；记作 \(\mathcal{Calk}(\mathrm{E})\) 。我们用 \(\pi\) 表示从 \(\mathcal{L}(\mathrm{E})\) 到 \(\mathcal{Calk}(\mathrm{E})\) 上的典范同态。由 III, p. 74 的推论 1，E 的自同态 \(u\) 是弗雷德霍姆自同态，当且仅当 \(\pi (u)\) 在代数 \(\mathcal{Calk}(\mathrm{E})\) 中可逆。

命题 2. --- 设 \(u \in \mathcal{L}(\mathrm{E})\) 满足 \(\|1_{\mathrm{E}} - \pi(u)\| < 1\) （在代数 \(\mathcal{Calk}(\mathrm{E})\) 中）。则 \(u\) 是 E 的里斯自同态。

设 \(r \geqslant 1\) 是使 \(\| 1_{\mathrm{E}} - \pi(u)\|^{r} < \frac{1}{2}\) 的整数。设 \(\mathrm{P} \in \mathbf{C}[\mathrm{X}]\) 是多项式 \(\frac{1 - (1 - \mathrm{X})^r}{\mathrm{X}}\) 。令 \(v = 1_{\mathrm{E}} - (1_{\mathrm{E}} - u)^r\) 。于是我们有 \(v = u\mathrm{P}(u)\) 与 \(\| 1_{\mathrm{E}} - \pi(v)\| < \frac{1}{2}\) 。由于 E 的自同态 \(u\) 与 \(\mathrm{P}(u)\) 交换，只需证明 \(v\) 是 E 的里斯自同态（III, p. 49, 命题 9）。

由于 \(\|1_{\mathrm{E}}-\pi(v)\|<1/2\) ，由空间 \(\mathcal{Calk}(\mathrm{E})\) 中商范数的定义，存在 E 的一个紧自同态 h 和 E 的一个自同态 w，使得 \(v=1_{E}+h+w\) 且 \(\|w\|<\frac{1}{2}\) 。由 I, p. 22 的推论 1，元素 \(1_{E}+w\) 是 E 的一个自同构。由于 h 是紧的，\(v=(1_{\mathrm{E}}+w)+h\) 是 E 的指标为 0 的弗雷德霍姆自同态（III, p. 74 的推论 2）。对每个整数 \(n\geqslant0\) ，用 \(N_{n}\) 表示 \(v^{n}\) 的核。为证明 v 是 E 的里斯自同态，只需证明存在一个整数 \(n \geqslant 0\) 使 \(N_{n} = N_{n+1}\) （III, p. 46, 定义 2 及 III, p. 46, 注记）。

我们用反证法，假设序列 \((\mathrm{N}_{n})\) 严格递增。对每个 \(n \in N\) ，用 \(p_{n}\) 表示从 E 到赋范空间 \(E/N_{n}\) 上的典范映射。设 c 是使 \(2\|w\| < c < 1\) 的实数。设 \(n \in N\) 。由于 \(N_{n+1}\) 异于 \(N_{n}\) ，存在一个元素 \(x_{n} \in N_{n+1}\) ，使得 \(\|p_{n}(x_{n})\| = c\) 且 \(\|x_{n}\| < 1\) （事实上，存在 \(y \in N_{n+1}/N_{n}\) 范数为 c，故对每个 \(\varepsilon > 0\) ，存在 \(x_{n} \in N_{n+1}\) 使得 \(p_{n}(x_{n}) = y\) 且 \(\|x_{n}\| \leqslant c + \varepsilon\) ）。

设 \(m\) 、\(n\) 是使 \(m > n \geqslant 0\) 的两个整数。我们有

\[
\begin{array}{r l} & {\| h (x _ {m}) - h (x _ {n}) \| = \| v (x _ {m} - x _ {n}) - (1 _ {\mathrm{E}} + w) (x _ {m} - x _ {n}) \|} \\ & {\qquad \geqslant \| (v - 1 _ {\mathrm{E}}) (x _ {m} - x _ {n}) \| - \| w (x _ {m} - x _ {n}) \|} \\ & {\qquad \geqslant \| (v - 1 _ {\mathrm{E}}) (x _ {m} - x _ {n}) \| - 2 \| w \|} \end{array}\tag{1}
\]

因为 \(\| x_m\| \leqslant 1\) 且 \(\| x_n\| \leqslant 1\) 。此外，注意到

\[
(v - 1 _ {\mathrm{E}}) (x _ {m} - x _ {n}) = v (x _ {m} - x _ {n}) + x _ {n} - x _ {m}.
\]

但由于 n \textless{} m 且 \(v(\mathrm{N}_{m+1}) \subset \mathrm{N}_{m}\) ，我们有 \(v(x_{m} - x_{n}) + x_{n} \in \mathrm{N}_{m}\) ，由此

\[
\left\| v (x _ {m} - x _ {n}) + x _ {n} - x _ {m} \right\| \geqslant \left\| p _ {m} (x _ {m}) \right\| = c.
\]

于是不等式 (1) 给出下界

\[
\left\| h \left(x _ {m}\right) - h \left(x _ {n}\right) \right\| \geqslant c - 2 \| w \| > 0.\tag{2}
\]

于是序列 \((h(x_{m}))_{m\in\mathbf{N}}\) 在 E 中没有聚点；由于 h 是紧的（III, p. 2 的注记 1），E 的单位球在 h 下的像是相对紧的，这就产生矛盾。

推论 1. --- 设 E 是分离局部凸空间，且 \(u \in \mathcal{L}(\mathrm{E})\) 。为使 u 是里斯自同态，必须且只需存在 \(\mathcal{L}(\mathrm{E})\) 的一个元素 v，它与 u 交换且使 \(1_{\mathrm{E}} - u \circ v\) 是紧的。

由于每个有限秩连续线性映射都是紧的，该条件是必要的（III, p. 47, 命题 6 d)）。我们来证明它是充分的。设 v 是 \(\mathcal{L}(\mathrm{E})\) 的一个元素，它与 u 交换，且使 E 的自同态 \(h = 1_{E} - u \circ v\) 是紧的。存在一个巴拿赫空间 G、一个连续线性映射 \(p: E \to G\) 和一个紧线性映射 \(q: G \to E\)，使得 \(h = q \circ p\) （III, p. 5, 命题 5）。G 的自同态 \(p \circ q\) 是紧的。由命题 2，线性映射 \(1_{G} - p \circ q\) 是 G 的里斯自同态。但此时 \(u \circ v = 1_{E} - q \circ p\) 是 E 的里斯自同态（III, p. 49, 命题 10, g)）。由于 u 与 v 交换，u 是 E 的里斯自同态（III, p. 49, 命题 9）。

推论 2. --- 设 E 是分离局部凸空间，u 是 E 的里斯自同态，\(h \in \mathcal{L}(\mathrm{E})\) 是紧线性映射。若 u 与 h 交换，则 \(u + h\) 是 E 的里斯自同态。

设 \(v\) 是 \(u\) 的典范拟逆。它与 \(u\) 交换，也与 \(h\) 交换（III, p. 47, 命题 6），从而与 \(u + h\) 交换。E 的自同态 \(1_{\mathrm{E}} - (u + h) \circ v\) 是有限秩连续线性映射 \(1_{\mathrm{E}} - u \circ v\) 与紧线性映射 \(-h \circ v\) 之和，故是紧的。由此得出 \(u + h\) 是 E 的里斯自同态（推论 1）。

命题 3. --- 设 E 是巴拿赫空间。设 u 是 E 的一个自同态，A 是它在 \(\mathcal{L}(\mathrm{E})\) 中的交换子。\(\pi(\mathrm{A})\) 在代数 \(\mathcal{Calk}(\mathrm{E})\) 中的闭包 B 是一个巴拿赫代数。为使 u 是 E 的里斯自同态，必须且只需 \(\pi(u)\) 在 B 中可逆。

由于 \(\pi(A)\) 是 \(\mathcal{Calk}(E)\) 的一个含单位元的子代数，其闭包 B 是一个巴拿赫代数。若 \(u\) 是 E 的里斯自同态，则元素 \(\pi(u)\) 在 \(\pi(A)\) 中有一个逆（推论 1），从而在 B 中有逆。反之，假设 \(\pi(u)\) 在 B 中有一个逆 \(s\) 。由 B 的定义，存在 A 的一个元素 \(v\) 使得

\[
\left\| s - \pi (v) \right\| <   \frac {1}{\left\| \pi (u) \right\|},
\]

由此 \(\|1_{E}-\pi(u\circ v)\|<1\) 。由命题 2，线性映射 \(u\circ v\) 是 E 的里斯自同态。由于 u 与 v 交换，u 也是（III, p. 49, 命题 9）。

\subsection{弗雷德里克·里斯理论}

设 E 是分离局部凸空间，h 是 E 的一个紧自同态。设 \(\lambda \in K\) 。自同态 \(\lambda h\) 是紧的，故 \(1_{E} - \lambda h\) 是 E 的里斯自同态（III, p. 75 的命题 2 的推论 2）；用 \(N_{\lambda}\) 与 \(I_{\lambda}\) 表示它的幂零空间与余幂零空间。由 III, p. 47 的命题 6，我们有下列性质：

\begin{enumerate}
\def\labelenumi{(\roman{enumi})}
\item
  E 的向量子空间 \(I_{\lambda}\) 与 \(N_{\lambda}\) 是闭的且在 h 下稳定；
\item
  局部凸空间 E 是 \(I_{\lambda}\) 与 \(N_{\lambda}\) 的拓扑直和；
\item
  映射 \(1_{\mathrm{E}} - \lambda h\) 通过限制诱导 \(\mathrm{I}_{\lambda}\) 的一个自同构；
\item
  向量空间 \(N_{\lambda}\) 有限维，且存在一个整数 \(n_{\lambda} \geqslant 0\) ，使得 \((1_{\mathrm{E}} - \lambda h)^{n_{\lambda}}(x) = 0\) （对一切 \(x \in N_{\lambda}\) ）。
\end{enumerate}

这些性质以唯一的方式确定空间 \(I_{\lambda}\) 与 \(N_{\lambda}\) （A, VIII, p. 26, 注记 2）。

引理 3. --- 设 X 是一个拓扑具有可数基的拓扑空间。X 的每个离散子集都是可数的。

设 \(\mathcal{B}\) 是 X 的拓扑的一个可数基，D 是 X 的一个离散子集。对 D 的每个元素 \(x\) ，存在元素 \(\mathrm{B}_x\) （属于 \(\mathcal{B}\) ），使得 \(\mathrm{B}_x \cap \mathrm{D} = \{x\}\) 。从 D 到 \(\mathcal{B}\) 的映射由 \(x \mapsto \mathrm{B}_x\) 定义，它是单射，由此引理得证。

定理 4. --- 设 E 是分离局部凸空间，h 是 E 的一个紧自同态。集 \(\Sigma\) 由那些 \(\lambda \in K\) （使 \(1_{\mathrm{E}} - \lambda h\) 不是 E 的自同构的）组成，它是域 K 的一个可数、闭且离散的子集。

集 \(\Sigma\) 也是由那些 \(\lambda\in\mathrm{K}\) （使 \(1_{\mathrm{E}}-\lambda h\) 不是单射的）所成的集，以及由那些 \(\lambda\in\mathrm{K}\) （使 \(1_{\mathrm{E}}-\lambda h\) 不是满射的）所成的集。

存在一个巴拿赫空间 G、一个连续线性映射 \(p: E \rightarrow G\) 和一个紧线性映射 \(q: G \rightarrow E\)，使得 \(h = q \circ p\) （III, p. 5, 命题 5）。G 的自同态 \(h' = p \circ q\) 是紧的，且 \(1_{E} - \lambda h\) 是 E 的自同构当且仅当 \(1_{G} - \lambda h'\) 是 G 的自同构（III, p. 49 的命题 10, e)）。因此我们只需在 E 是巴拿赫空间的情形证明该定理，以下即作此假设。

设 \(\lambda \in K\) 。由于里斯自同态的指标为零，下列说法等价：映射 \(1_{\mathrm{E}} - \lambda h\) 是 E 的自同构，或它是单射的，或它是满射的。

集 \(\Sigma\) 在 K 中是闭的，因为 E 的自同构所成的集在 \(\mathcal{L}(\mathrm{E})\) 中是开的。设 \(\lambda\) 是 \(\Sigma\) 的一个元素。我们有 \(\lambda \neq 0\) 。用 N 表示 h 的幂零空间，I 表示其余幂零空间，并用 \(h_{I}\) 与 \(h_{N}\) 表示 h 在这些子空间上诱导的 I 与 N 的自同态。则 \(1_{I} - \lambda h_{I}\) 是 I 的一个自同构，且存在 \(\lambda\) 在 K 中的一个邻域 U，使得 \(1_{I} - \mu h_{I}\) 是 I 的自同构（对一切 \(\mu \in U\) ）。N 的自同态 \(1_{N} - \lambda h_{N}\) 是幂零的；对一切 \(\mu \neq \lambda\) ，我们有

\[
1 _ {\mathrm{N}} - \mu h _ {\mathrm{N}} = \frac {\lambda - \mu}{\lambda} \Big (1 _ {\mathrm{N}} + \frac {\mu}{\lambda - \mu} (1 _ {\mathrm{N}} - \lambda h _ {\mathrm{N}}) \Big),
\]

故 \(1_{N} - \mu h_{N}\) 是 N 的一个自同构。于是集 \(U \cap \Sigma\) 缩减为单个元素 \(\lambda\) ，从而集 \(\Sigma\) 是离散的。由引理 3 它是可数的。

\subsection{弗雷德霍姆择一定理}

命题 4（弗雷德霍姆择一定理）. --- 设 E 是巴拿赫空间，h 是 E 的一个紧自同态。设 \(\lambda\) 是 K-\{0\} 的一个元素。设 F 是 \(1_{\mathrm{E}} - \lambda h\) 的核，G 是 \(1_{\mathrm{E}'} - \lambda^{t}h\) 的核。

\begin{enumerate}
\def\labelenumi{\alph{enumi})}
\item
  空间 F 与 G 有相同的有限维数 n ≥ 0；
\item
  对 \(y \in \mathrm{E}\) ，存在 \(x \in \mathrm{E}\) 使 \(x - \lambda h(x) = y\) ，当且仅当 \(\langle y, \ell \rangle = 0\) （对一切 \(\ell \in \mathrm{G}\) ）；
\item
  对 \(\ell \in \mathrm{E}'\) ，存在 \(f \in \mathrm{E}'\) 使 \(f - \lambda^t h(f) = \ell\) ，当且仅当 \(\langle x, f \rangle = 0\) （对一切 \(x \in \mathrm{F}\) ）。
\end{enumerate}

特别地，下列条件中恰有一个成立：

\begin{enumerate}
\def\labelenumi{(\roman{enumi})}
\item
  空间 F 非零；
\item
  对每个 \(y \in \mathrm{E}\) ，存在 \(x \in \mathrm{E}\) 使 \(x - \lambda h(x) = y\) 。
\end{enumerate}

通过把 h 换成 \(\lambda h\) ，我们化归到 \(\lambda = 1\) 的情形。

自同态 \(^{t}h\) 由 III, p. 9 的推论 1 是紧的。于是 E 的自同态 \(w = 1_{\mathrm{E}} - h\) 与 \(w' = 1_{\mathrm{E}'} - ^{t}h\) （\(\mathrm{E}'\) 的）都是里斯自同态（III, p. 75 的命题 2 的推论 2）。特别地，它们的核 F 与 G 是有限维的。由于 \(w' = ^{t}w\) 且 w 的指标为零，由 III, p. 43 的公式 (3)，F 的维数等于 G 的维数。

对 E 的元素 \(y\) ，方程 \(x - h(x) = y\) 有解 \(x \in \mathrm{E}\) 当且仅当 \(x\) 属于 \(w\) 的像。由于 \(w\) 的余核的对偶与 \(\operatorname{Ker}(^t w)\) 的核 G 等同（EVT, IV, p. 27, 命题 2），这件事成立当且仅当 \(\langle y, \ell \rangle = 0\) （对一切 \(\ell \in \mathrm{G}\) ）。

对 E' 的元素 \(\ell\) ，方程 \(f - {}^t h(f) = \ell\) 有解 \(f \in \mathrm{E}'\) 当且仅当 \(f\) 属于 \(w'\) 的像。由于 \({}^t w\) 的像是 \(w\) 的核 F 的正交（EVT, IV, p. 27, 命题 2），这件事成立当且仅当 \(\langle x, f \rangle = 0\) （对一切 \(x \in \mathrm{F}\) ）。

例. --- 设 X 是紧空间，\(\mu\) 是 X 上的一个测度（依 K 等于 R 或 C 而为实或复），N: \(X \times X \to K\) 是一个连续函数。

用 E 表示巴拿赫空间 \(\mathcal{C}(\mathrm{X};\mathrm{K})\) 。假设给定了函数 \(a\in\mathrm{E}\) 与 \(\lambda\in\mathrm{K}-\{0\}\) 。我们研究积分方程的解 \(f\in\mathrm{E}\) 的存在性与唯一性问题

\[
f (x) - \lambda \int_ {\mathrm{X}} \mathrm{N} (x, y) f (y) d \mu (y) = a (x) \quad (x \in \mathrm{X}).\tag{3}
\]

为此，引入集 \(F_{\lambda}\) （由相应的齐次方程的解 \(f \in E\) 组成）

\[
f (x) - \lambda \int_ {\mathrm{X}} \mathrm{N} (x, y) f (y) d \mu (y) = 0 \qquad (x \in \mathrm{X}),\tag{4}
\]

以及集 \(G_{\lambda}\) （由 (4) 的转置齐次方程的解 \(g \in E\) 组成），即

\[
g (y) - \lambda \int_ {\mathrm{X}} \mathrm{N} (x, y) g (x) d \mu (x) = 0 \qquad (y \in \mathrm{X}).\tag{5}
\]

若 (3) 的解集非空，则它是 E 的一个以 \(F_{\lambda}\) 为方向的仿射子空间。

对 \(f\in \mathrm{E}\) 与 \(x\in \mathrm{X}\) ，令

\[
h (f) (x) = \int_ {\mathrm{X}} \mathrm{N} (x, y) f (y) d \mu (y);
\]

如此定义的映射 \(f \mapsto h(f)\) 是巴拿赫空间 E 的一个紧自同态（III, p. 26 的命题 2 的推论）。E 的对偶 E' 由 X 上的测度组成（依 K 等于 R 或 C 而为实或复）（INT, III, p. 47, § 1, 第 3 目, 定义 2）。h 的转置 \(^{t}h\) 是 E' 的自同态，由关系式 \(\langle f, ^{t}h(m)\rangle = \langle h(f), m\rangle\) （对 \(f \in E\) 与 \(m \in E'\) ）刻画，即

\[
\begin{array}{r l} \int_ {\mathrm{X}} f (x) d (^ {t} h (m)) (x) & = \int_ {\mathrm{X}} h (f) (x) d m (x) \\ & = \int_ {\mathrm{X}} \Big (\int_ {\mathrm{X}} \mathrm{N} (x, y) f (y) d \mu (y) \Big) d m (x) \\ & = \int_ {\mathrm{X}} \Big (\int_ {\mathrm{X}} \mathrm{N} (x, y) d m (x) \Big) f (y) d \mu (y) \end{array}\tag{6}
\]

（对 \(m \in E'\) 与 \(f \in E\) ；INT, III, p. 84, § 4, 第 1 目, 定理 2）。

引理 4. --- 从 \(G_{\lambda}\) 到 \(E'\) 中的线性映射（由 \(g \mapsto g \cdot \mu\) 定义）是单射的，其像就是 \(1_{E'} - \lambda^{t}h\) 的核。

若 \(g \in G_{\lambda}\) 满足 \(g \cdot \mu = 0\) ，则公式 (5) 蕴含 \(g = 0\) ，故映射 \(g \mapsto g \cdot \mu\) 是单射的。

设 \(m \in \mathrm{E}'\) 是属于 \(1_{\mathrm{E}'} - \lambda^t h\) 的核的一个测度。由关系式 (6)，我们有

\[
\int_ {\mathrm{X}} f (x) d m (x) = \lambda \int_ {\mathrm{X}} \left(\int_ {\mathrm{X}} \mathrm{N} (x, y) d m (x)\right) f (y) d \mu (y)\tag{7}
\]

（对每个连续函数 \(f: \mathrm{X} \to \mathrm{K}\) ）。于是测度 \(m\) 等于测度 \(g \cdot \mu\) ，其中 \(g\) 是由下式定义的从 \(\mathrm{X}\) 到 \(\mathrm{K}\) 中的连续函数

\[
g (y) = \lambda \int_ {\mathrm{X}} \mathrm{N} (x, y) d m (x)\tag{8}
\]

（对一切 \(y \in X\) ）。于是公式 (8) 蕴含函数 g 属于 \(G_{\lambda}\) 。

反之，设 \(g: X \to K\) 是属于 \(G_{\lambda}\) 的一个连续函数，令 \(m = g \cdot \mu\) 。于是对每个连续函数 \(f: X \to K\) ，我们有

\[
\begin{array}{l} \int_ {\mathrm{X}} f (y) d m (y) = \int_ {\mathrm{X}} g (y) f (y) d \mu (y) \\ \qquad = \lambda \int_ {\mathrm{X}} \Big (\int_ {\mathrm{X}} \mathrm{N} (x, y) g (x) d \mu (x) \Big) f (y) d \mu (y) \\ \qquad = \lambda \int_ {\mathrm{X}} f (x) d m (x), \end{array}
\]

从而 \({}^{t}h(m)=\lambda m\) 。这就证明了引理 4。

把 III, p. 78 的命题 4 与定理 4 应用于自同态 h，我们便得到下列命题。

定理 5. --- 设 \(\lambda\in K\) 。

\begin{enumerate}
\def\labelenumi{\alph{enumi})}
\item
  集 \(F_{\lambda}\) 与 \(G_{\lambda}\) 是 \(\mathcal{C}(X;K)\) 的有限维向量子空间，且它们的维数相等；
\item
  为使方程 (3) 至少有一个解 \(f \in \mathcal{C}(\mathrm{X}, \mathrm{K})\) ，必须且只需有 \(\int_{\mathrm{X}} a(x) g(x) d\mu(x) = 0\) （对每个函数 \(g \in G_{\lambda}\) ）。此时 (3) 的解集是 \(\mathcal{C}(\mathrm{X}; \mathrm{K})\) 的一个以 \(F_{\lambda}\) 为方向的仿射子空间；
\item
  下列互相排斥的条件之一成立：
\end{enumerate}

\begin{enumerate}
\def\labelenumi{(\roman{enumi})}
\item
  对任何函数 \(a \in \mathcal{C}(\mathrm{X};\mathrm{K})\) ，积分方程 (3) 存在唯一的解 \(f \in \mathcal{C}(\mathrm{X};\mathrm{K})\) ；
\item
  齐次方程 (4) 有一个非零解 \(f \in \mathcal{C}(\mathrm{X};\mathrm{K})\) 。
\end{enumerate}

推论. --- 由 \(\lambda\in\mathrm{K}\) （使空间 \(\mathrm{F}_{\lambda}\) 非零的）组成的集是 K 的一个可数、闭且离散的子集。

\section{巴拿赫空间自同态的谱性质}

除另有说明外，本段所考虑的向量空间都是 C 上的向量空间。完备可赋范空间 E 的自同态的谱是相对于含单位元代数 \(\mathcal{L}(\mathrm{E})\) 的谱（参见 I, p. 127 的第 7 目）。

\subsection{谱的孤立点与敏感点}

设 E 是完备可赋范空间，u 是 E 的一个自同态。设 \(\lambda\) 是 u 的谱的一个孤立点。回忆我们用 \(e_{\lambda}(u)\) 表示与 u 及 u 的谱的既闭又开的子集 \(\{\lambda\}\) 相伴的谱投影（I, p. 131 的第 3 目）。自同态 \(u - \lambda1_{E}\) 诱导 \(\operatorname{Ker}(e_{\lambda}(u))\) 的一个自同构以及 \(\operatorname{Im}(e_{\lambda}(u))\) 的一个拟幂零自同态（同前）。

定义 1. --- 我们说 \(\lambda\) 对 \(u\) 有有限谱重数，如果谱投影 \(e_{\lambda}(u)\) 有限秩；此时，整数 \(m_{\lambda}(u) = \dim (\mathrm{Im}(e_{\lambda}(u)))\) 称为 \(\lambda\) 对 \(u\) 的谱重数。

定义 2. --- 我们把 Sp(u) 的有限谱重数的孤立点称为 u 的谱的敏感点。Sp(u) 的敏感点所成的集称为 u 的敏感谱，记作 Sp \(_{s}\) (u)。\(^{(2)}\) 我们把 u 的本质谱记作 Sp \(_{e}\) (u)，即补集 \(\mathrm{Sp}(u) - \mathrm{Sp}_{\mathrm{s}}(u)\)。

由于集 \(\mathrm{Sp}_{\mathrm{s}}(u)\) 由 \(\mathrm{Sp}(u)\) 的孤立点组成，它在 \(\mathrm{Sp}(u)\) 中是开的、离散的且可数的（III, p. 78, 引理 3）；它还是有界的。集 \(\mathrm{Sp}_{\mathrm{e}}(u)\) 在 \(\mathrm{Sp}(u)\) 中是闭的，故是紧的。

由 III, p. 54 的命题 14，\(u\) 的谱的敏感点是那些复数 \(\lambda\) （使 \(u - \lambda 1_{\mathrm{E}}\) 为 E 的里斯自同态且不是 E 的自同构的）。

对每个复数 \(\lambda\) ，用 \(\mathrm{N}_{\lambda}(u)\) 与 \(\mathrm{I}_{\lambda}(u)\) 表示 E 的自同态 \(u - \lambda 1_{\mathrm{E}}\) 的幂零空间与余幂零空间。由于 E 是弗雷歇空间，\(u - \lambda 1_{\mathrm{E}}\) 是里斯自同态当且仅当 \(\mathrm{N}_{\lambda}(u)\) 有限维且 \(\mathrm{I}_{\lambda}(u)\) 有有限余维数（III, p. 53 的命题 13）；在此情形，\(u - \lambda 1_{\mathrm{E}}\) 可逆当且仅当 \(\mathrm{N}_{\lambda}(u)\) 缩减为 0（III, p. 47, 命题 6）。于是 \(u\) 的谱的敏感点是复数 \(\lambda\) （使 \(\mathrm{N}_{\lambda}(u)\) 为非零有限维且 \(\mathrm{I}_{\lambda}(u)\) 在 E 中有有限余维数的）（III, p. 54, 命题 14）。此时谱重数 \(m_{\lambda}(u)\) （\(\lambda\) 的）就是 \(\mathrm{N}_{\lambda}(u)\) 的维数，且自同态 \(u - \lambda 1_{\mathrm{E}}\) 通过限制诱导 \(\mathrm{I}_{\lambda}(u)\) 的一个自同构与 \(\mathrm{N}_{\lambda}(u)\) 的一个幂零自同态。由于 \(\mathrm{N}_{\lambda}(u)\) 不缩减为 0，由此得出 \(\lambda\) 是 \(u\) 的一个特征值。最小的整数 \(n \geqslant 1\) （使 \(\operatorname{Ker}((u - \lambda 1_{\mathrm{E}})^n)\) 等于 \(\mathrm{N}_{\lambda}(u)\) ）是 \(u\) 的预解式在点 \(\lambda\) 处极点的阶（I, p. 131 的第 3 目），它以 \(m_{\lambda}(u)\) 为上界。

特别地，我们由此证明了：

命题 1. --- u 的谱的敏感点是有限谱重数的特征值。

\(u\) 的本质谱由那些复数 \(\lambda\) （使 \(u - \lambda 1_{\mathrm{E}}\) 不是 E 的里斯自同态的）组成。特别地，若 \(u\) 是紧的，则每个 \(\lambda \in \mathrm{Sp}(u) - \{0\}\) 都是谱的敏感点（III, p. 77 的推论 2）。

命题 2. --- 设 E 是完备可赋范空间，u 是 E 的一个自同态。设 H 是 \(\mathrm{Sp}_{\mathrm{s}}(u)\) 的一个有限子集。集 H 在 \(\mathrm{Sp}(u)\) 中既开又闭。谱投影 \(e_{\mathrm{H}}(u)\) 有限秩，其像为 \(\sum_{\lambda \in \mathrm{H}}\mathrm{N}_{\lambda}(u)\) ，核为 \(\bigcap_{\lambda \in \mathrm{H}}\mathrm{I}_{\lambda}(u)\) 。

这由 I, p. 131 的第 3 目得出，因为对每个 \(\lambda \in \mathrm{Sp}_{\mathrm{s}}(u)\) ，子空间 \(\mathrm{N}_{\lambda}(u)\) 与 \(\mathrm{I}_{\lambda}(u)\) 分别是谱投影 \(e_{\lambda}(u)\) 的像与核。

推论 1. --- 设 V 是 \(\mathrm{Sp_e}(u)\) 的一个邻域。

\begin{enumerate}
\def\labelenumi{\alph{enumi})}
\item
  存在 E 作为拓扑直和 F ⊕ G 的分解，使得 F 有限维，F 与 G 在每个与 u 交换的自同态下稳定，且 u 诱导的 G 的自同态的谱包含于 V。
\item
  假设 V 非空。存在元素 v （属于 u 在 \(\mathcal{L}(\mathrm{E})\) 中的双交换子），其谱包含于 V 且 \(v - u\) 有限秩。
\end{enumerate}

令 \(\mathrm{H} = \mathrm{Sp}(u) \cap (\mathbf{C} - \mathring{\mathrm{V}})\) 。集 H 包含于 \(\mathrm{Sp}_{\mathrm{s}}(u)\) ，故是离散的；由于它在紧空间 \(\mathrm{Sp}(u)\) 中是闭的，它是有限的。可取谱投影 \(e_{\mathrm{H}}(u)\) 的像与核作为 F 与 G（命题 2）。

假设 V 非空。设 \(\mu\in\mathrm{V}\) ，用 \(v\) 表示 E 的这样一个自同态：它在 F 上与比率为 \(\mu\) 的位似重合，在 G 上与 \(u\) 重合。其谱包含于 \(\{\mu\}\cup(\mathrm{Sp}(u)-\mathrm{H})\) ，从而包含于 V，且 \(v-u\) 有限秩，因为 F 有限维。

对 C 的每个紧子集 S，我们用 \(\mathcal{O}(S)\) 表示在 S 的一个邻域中定义、取值于 C 的全纯函数芽所成的代数（I, p. 49, §4, 第 1 目）。

推论 2. --- 设 \(f \in \mathcal{O}(\mathrm{Sp}(u))\) ，\(\mu\) 是一个复数。设 H 是 \(\lambda \in \mathrm{Sp}(u)\) （满足 \(f(\lambda) = \mu\) ）所成的集。则 \(\mu\) 是 \(f(u)\) 的谱的敏感点，当且仅当 H 有限、非空且包含于 \(u\) 的敏感谱。在这些条件下，谱投影 \(e_{\mu}(f(u))\) 等于投影 \(e_{\mathrm{H}}(u)\) （与 \(u\) 及 H 相伴的）。我们有

\[
\mathrm{N} _ {\mu} (f (u)) = \bigoplus_ {\lambda \in \mathrm{H}} \mathrm{N} _ {\lambda} (u), \quad \mathrm{I} _ {\mu} (f (u)) = \bigcap_ {\lambda \in \mathrm{H}} \mathrm{I} _ {\lambda} (u)
\]

而 \(\mu\) 对 \(f(u)\) 的谱重数是 H 的元素的谱重数之和，即

\[
m _ {\mu} (f (u)) = \sum_ {\lambda \in \mathrm{H}} m _ {\lambda} (u).
\]

复数 \(\mu\) 属于 \(f(\mathrm{Sp}(u))\) ，从而属于 \(f(u)\) 的谱（I, p. 75 的命题 8），当且仅当 H 非空。此时我们有 \(e_{\mu}(f(u)) = e_{\mathrm{H}}(u)\) （I, p. 127 的第 1 目）。其余断言由命题 2 得出。

例. --- 1) 设 E 是有限维向量空间，u 是 E 的一个自同态。u 的谱有限且与 \(\mathrm{Sp}_{\mathrm{s}}(u)\) 重合。其元素是 u 的特征多项式 \(\chi_{u}\) 的根；它们就是 u 的特征值。谱重数 \(m_{\lambda}(u)\) （一个元素 \(\lambda \in \mathrm{Sp}(u)\) 的）是 \(\lambda\) 作为多项式 \(\chi_{u}\) 的根的重数（A, VII, p. 36, 推论）。由上面的推论 2，我们有

\[
\chi_ {f (u)} (\mathrm{T}) = \prod_ {\lambda \in \operatorname{Sp} (u)} (\mathrm{T} - f (\lambda)) ^ {m _ {\lambda}}
\]

（对一切 \(f \in \mathcal{O}(\mathrm{Sp}(u))\) ）。这推广了 A, VII, p. 36 的命题 10。

\begin{enumerate}
\def\labelenumi{\arabic{enumi})}
\setcounter{enumi}{1}
\tightlist
\item
  设 X 是 C 的一个紧子集，\(\mathcal{C}(X)\) 是 X 上取复值的连续函数所成的巴拿赫空间。用 u 表示 \(\mathcal{C}(X)\) 的这样一个自同态：它把 \(f \in \mathcal{C}(X)\) 映为函数 \(z \mapsto zf(z)\) （在 \(\mathcal{C}(X)\) 中） 。u 的谱等于 X，其敏感点是 X 的孤立点，且它们的谱重数等于 1。
\end{enumerate}

\subsection{谱的一个分划}

用 \(\overline{Z}\) 表示子集 \(Z \cup \{-\infty, +\infty\}\) （\(\overline{R}\) 中的）。若 u 是一个核或余核有限维的线性映射，则称 u 的指标为由下式定义的元素 \(\text{ind}(u)\) （\(\overline{Z}\) 中的）

\[
\operatorname{ind} (u) = \dim \operatorname{Coker} (u) - \dim \operatorname{Ker} (u)
\]

（参见 III, p. 67 的第 6 目）。

定义 3. --- 设 E 是完备可赋范空间，u 是 E 的一个自同态。对每个 \(n \in \overline{Z}\) ，用 \(\mathrm{Sp}_{n}(u)\) 表示这样的复数 \(\lambda \in \mathrm{Sp}_{\mathrm{e}}(u)\) 所成的集：使 \(u - \lambda1_{E}\) 是严格态射，其核或余核有限维，且其指标为 \(n\) 。我们用 \(\mathrm{Sp}_{\omega}(u)\) 表示在 \(\mathrm{Sp}_{\mathrm{e}}(u)\) 中子集 \(\mathrm{Sp}_{n}(u)\) （\(n \in \overline{Z}\) ）之并的补集。

诸集 \(\mathrm{Sp}_{\mathrm{s}}(u)\) 、\(\mathrm{Sp}_{n}(u)\) （\(n \in \overline{Z}\) ）与 \(\mathrm{Sp}_{\omega}(u)\) 构成 u 的谱的一个分划。

E 的每个余核有限维的自同态都是严格的（III, p. 52, 引理 6）。E 的弗雷德霍姆自同态就是 E 的核与余核均有限维的自同态（III, p. 52, 命题 11）。集 \(\mathbf{C} - \mathrm{Sp}_{\mathrm{e}}(u)\) 由那些 \(\lambda \in \mathbf{C}\) （使 \(u - \lambda 1_{\mathrm{E}}\) 是 E 的里斯自同态的）组成，而这样的自同态是 E 的指标为 0 的弗雷德霍姆自同态。

于是，对 \(\lambda \in \mathbf{C}\) 及 \(n\in \mathbf{Z} - \{0\}\) ，我们有：

\(\lambda \in \mathbf{C} - \mathrm{Sp}(u) \iff u - \lambda 1_{\mathrm{E}}\) 是自同构；

\(\lambda \in \mathrm{Sp}_{\mathrm{s}}(u) \quad \Longleftrightarrow \quad u - \lambda 1_{\mathrm{E}} \text{ 是里斯自同态，但不是自同构；}\)

\(\lambda \in \mathrm{Sp}_{0}(u) \quad \Longleftrightarrow \quad u - \lambda 1_{\mathrm{E}}\) 是 E 的指标为 0 的弗雷德霍姆自同态但不是 E 的里斯自同态；

\(\lambda \in \mathrm{Sp}_n(u) \quad \Longleftrightarrow \quad u - \lambda 1_{\mathrm{E}}\) 是 E 的指标为 \(n\) 的弗雷德霍姆自同态（\(n \neq 0\)）；

\[
\lambda \in \mathrm{Sp} _ {- \infty} (u) \quad \Longleftrightarrow \quad u - \lambda 1 _ {\mathrm{E}} \text{ 是严格态射，其核无限维且其余核有限维；}
\]

\(\lambda \in \mathrm{Sp}_{+\infty}(u) \quad \Longleftrightarrow \quad u - \lambda 1_{\mathrm{E}}\) 是严格态射，其核有限维且其余核无限维；

\(\lambda \in \mathrm{Sp}_{\omega}(u) \quad \Longleftrightarrow \quad \text{要么 } u - \lambda 1_{\mathrm{E}} \text{ 不是严格态射，要么其核与余核都是无限维。}\)

注记. --- 用 \(\pi\) 表示从 \(\mathcal{L}(\mathrm{E})\) 到 E 的卡尔金代数 \(\mathcal{C}alk(\mathrm{E})\) 上的典范同态（参见 III, p. 75 的第 4 目）。\(\pi(u)\) 相对于代数 \(\mathcal{C}alk(\mathrm{E})\) 的谱有时称为 \(u\) 的稳定谱。其元素是那些复数 \(\lambda\) （使 \(u - \lambda 1_{\mathrm{E}}\) 不是 E 的弗雷德霍姆自同态的）（III, p. 73 的定理 3 的推论 1）。因此它等于 \(\mathrm{Sp}_{\omega}(u) \cup \mathrm{Sp}_{-\infty}(u) \cup \mathrm{Sp}_{+\infty}(u)\) 。

设 A 是 E 的与 u 交换的自同态所成的集。\(\pi(\mathrm{A})\) 在 \(\mathcal{C}alk(\mathrm{E})\) 中的闭包 B 是一个完备可赋范代数，且 \(\pi(u)\) 相对于 B 的谱就是 u 的本质谱 \(\mathrm{Sp}_{\mathrm{e}}(u)\) （III, p. 77, 命题 3）。把 I, p. 28 的命题 6 的推论应用于 \(\mathcal{C}alk(\mathrm{E})\) 的子代数 B，便知集 \(\mathrm{Sp}_{\mathrm{e}}(u)\) 是 \(\mathrm{Sp}_{\omega}(u) \cup \mathrm{Sp}_{-\infty}(u) \cup \mathrm{Sp}_{+\infty}(u)\) 与其补集的某些有界连通分支的并。

定理 1. --- 设 E 是完备可赋范空间，u 是 E 的一个自同态。

\begin{enumerate}
\def\labelenumi{\alph{enumi})}
\item
  集 \(\mathrm{Sp}_{\omega}(u)\) 是紧的。若 E 无限维，则它非空；
\item
  设 \(n \in \overline{Z}\) 。集 \(\mathrm{Sp}_{n}(u)\) 是 \(\mathbf{C} - \mathrm{Sp}_{\omega}(u)\) 的一族有界连通分支的并。它在 C 中是开的，且它在 C 中的边界包含于 \(\mathrm{Sp}_{\omega}(u)\) 。
\end{enumerate}

集 \(\mathbf{C} - \mathrm{Sp}_{\omega}(u)\) 由那些复数 \(\lambda \in \mathbf{C}\) （使 \(u - \lambda 1_{\mathrm{E}}\) 为核或余核有限维的严格态射的）组成。由 III, p. 67 的命题 11 与 III, p. 70 的命题 13，它是开的。于是集 \(\mathrm{Sp}_{\omega}(u)\) 是闭的。由于它有界，它是紧的。

我们来证明 \(b\) 。设 \(n \in \overline{\mathbf{Z}}\) 。集 \(\mathrm{Sp}_n(u)\) 包含于 \(\mathbf{C} - \mathrm{Sp}_{\omega}(u)\) 。设 U 是 \(\mathbf{C} - \mathrm{Sp}_{\omega}(u)\) 的一个与 \(\mathrm{Sp}_n(u)\) 相交的连通分支。由于映射 \(\lambda \mapsto \operatorname{ind}(u - \lambda 1_{\mathrm{E}})\) （从 \(\mathbf{C} - \mathrm{Sp}_{\omega}(u)\) 到 \(\overline{\mathbf{Z}}\) 中）是局部常数的（III, p. 68 的命题 12 的推论 1 及 III, p. 70 的命题 13 的推论 1），\(u - \lambda 1_{\mathrm{E}}\) 的指标等于 \(n\) （对一切 \(\lambda \in U\) ）。若 \(n \neq 0\) ，这就蕴含 U 包含于 \(\mathrm{Sp}_n(u)\) 。若 \(n = 0\) ，注意到集 U 是 \(\mathbf{C} - (\mathrm{Sp}_{\omega}(u) \cup \mathrm{Sp}_{-\infty}(u) \cup \mathrm{Sp}_{+\infty}(u))\) 的一个连通分支。由于 U 与 \(\mathrm{Sp}_0(u)\) 相交，从而与 \(\mathrm{Sp}_{\mathrm{e}}(u)\) 相交，由注记 2 得出集 U 包含于 \(\mathrm{Sp}_{\mathrm{e}}(u)\) ，从而包含于 \(\mathrm{Sp}_0(u)\) 。在所有情形，我们都得出结论：\(\mathrm{Sp}_n(u)\) 是 \(\mathbf{C} - \mathrm{Sp}_{\omega}(u)\) 的与 \(\mathrm{Sp}_n(u)\) 相交的那些连通分支的并。由于集 \(\mathrm{Sp}(u)\) 有界，这些连通分支必然有界。于是 \(\mathrm{Sp}_n(u)\) 在 \(\mathbf{C}\) 中是开的，且其边界包含于 \(\mathrm{Sp}_{\omega}(u)\) 。这就证明了 \(b\) 。

最后假设集 \(\mathrm{Sp}_{\omega}(u)\) 是空的。由 b)，诸集 \(\mathrm{Sp}_{n}(u)\) （\(n \in \overline{Z}\) ）此时也都是空的。于是我们有 \(\mathrm{Sp}(u) = \mathrm{Sp}_{\mathrm{s}}(u)\) 。因而 u 的谱是离散且紧的，故有限；又由于其所有点都有有限谱重数，向量空间 E 是有限维的（III, p. 83, 命题 2）。这就完成了 a) 的证明。

推论. --- a) 设 \(\Omega\) 是 \(\mathbf{C} - \mathrm{Sp}_{\omega}(u)\) 的无界连通分支。我们有 \(\Omega \cap \mathrm{Sp}(u) \subset \mathrm{Sp}_{\mathrm{s}}(u)\) ；

\begin{enumerate}
\def\labelenumi{\alph{enumi})}
\setcounter{enumi}{1}
\tightlist
\item
  每个附着于 \(\mathrm{Sp}_{\mathrm{s}}(u)\) 但不属于 \(\mathrm{Sp}_{\mathrm{s}}(u)\) 的点都属于 \(\mathrm{Sp}_{\omega}(u)\) 。
\end{enumerate}

断言 \(a\) 是定理 1 的断言 \(b\) 的直接推论。设 \(\lambda\) 是附着于 \(\mathrm{Sp}_{\mathrm{s}}(u)\) 但不属于 \(\mathrm{Sp}_{\mathrm{s}}(u)\) 的一点。它属于 \(u\) 的谱，因为后者是闭的。它不属于任何集 \(\mathrm{Sp}_{n}(u)\) （\(n \in \overline{\mathbf{Z}}\) ），因为这些集是开的（同前）且与 \(\mathrm{Sp}_{\mathrm{s}}(u)\) 不相交。于是我们有 \(\lambda \in \mathrm{Sp}_{\omega}(u)\) ，由此得 \(b\) 。

命题 3. --- 设 E 和 F 是完备可赋范空间，\(u\colon E\to F\) 与 \(v\colon F\to E\) 是连续线性映射。

\begin{enumerate}
\def\labelenumi{\alph{enumi})}
\item
  在 \(\mathbf{C} - \{0\}\) 上，集 \(\mathrm{Sp}(v \circ u)\) 与 \(\mathrm{Sp}(u \circ v)\) （相应地，\(\mathrm{Sp}_{\mathrm{s}}(v \circ u)\) 与 \(\mathrm{Sp}_{\mathrm{s}}(u \circ v)\) ，相应地，\(\mathrm{Sp}_{n}(v \circ u)\) 与 \(\mathrm{Sp}_{n}(u \circ v)\) （\(n \in \overline{\mathbf{Z}}\) ），相应地，\(\mathrm{Sp}_{\omega}(v \circ u)\) 与 \(\mathrm{Sp}_{\omega}(u \circ v)\) ）的迹相等；
\item
  设 \(\lambda\) 是 \(\mathrm{Sp}_{\mathrm{s}}(v \circ u)\) 中异于 0 的一个元素。\(\lambda\) 对 \(v \circ u\) 与对 \(u \circ v\) 的谱重数相等。
\end{enumerate}

设 \(\mu\) 是一个非零复数，且 \(n \in \overline{\mathbf{Z}}\) 。为使 \(\mu 1_{\mathrm{E}} - v \circ u\) 是一个自同构（相应地，一个里斯自同态，相应地，一个核或余核有限维且指标为 \(n\) 的严格态射），必须且只需 \(\mu 1_{\mathrm{F}} - u \circ v\) 也是（III, p. 49, 命题 10）。断言 \(a\) 于是由定义得出。

设 \(\lambda\) 是 \(\mathrm{Sp}_{\mathrm{s}}(v \circ u)\) 中异于 0 的一点。我们有

\[
\dim \operatorname{Ker} ((\lambda 1 _ {\mathrm{E}} - v \circ u) ^ {n}) = \dim \operatorname{Ker} ((\lambda 1 _ {\mathrm{F}} - u \circ v) ^ {n})
\]

（对一切 \(n \geqslant 0\) ；同前），故 \(\lambda\) 对 \(v \circ u\) 与 \(u \circ v\) 的谱重数相等。

\subsection{自同态的转置的谱}

命题 4. --- 设 E 是完备可赋范空间，E' 是 E 的对偶空间，u 是 E 的一个自同态。

\begin{enumerate}
\def\labelenumi{\alph{enumi})}
\item
  有 \(\mathrm{Sp}_{\mathrm{s}}(u) = \mathrm{Sp}_{\mathrm{s}}(^{t}u)\) ，\(\mathrm{Sp}_n(u) = \mathrm{Sp}_{-n}(^t u)\) （对一切 \(n \in \overline{\mathbf{Z}}\) ），且 \(\mathrm{Sp}_{\omega}(u) = \mathrm{Sp}_{\omega}(^t u)\) ；
\item
  \(\mathrm{Sp}_{\mathrm{s}}(u)\) 的每个点对 u 与 \({}^{t}u\) 有相同的谱重数。
\end{enumerate}

I, p. 131 的命题 3 的断言 \(a\) 表明 \(\mathrm{Sp}(u) = \mathrm{Sp}({}^{t}u)\) ，且同处的断言 \(c\) 蕴含 \(\mathrm{Sp}_{\mathrm{s}}(u) = \mathrm{Sp}_{\mathrm{s}}({}^{t}u)\) 以及断言 \(b\) 成立。

对每个 \(\lambda \in \mathbf{C}\) ，III, p. 69 的引理 4 蕴含：\(u - \lambda 1_{\mathrm{E}}\) 是严格态射当且仅当 \({}^{t}u - \lambda 1_{\mathrm{E}'}\) 是严格态射，且

\[
\dim \operatorname{Coker} (^ {t} (u - \lambda 1 _ {\mathrm{E}})) = \dim \operatorname{Ker} (u - \lambda 1 _ {\mathrm{E}})
\]

\[
\dim \operatorname{Ker} (^ {t} (u - \lambda 1 _ {\mathrm{E}})) = \dim \operatorname{Coker} (u - \lambda 1 _ {\mathrm{E}})
\]

（在 \(\overline{Z}\) 中）。由此并注意到谱的各部分的定义（III, p. 85 的定义 3），便得出断言 a)。

\subsection{借紧算子的扰动}

定理 2. --- 设 E 是完备可赋范空间，u 是 E 的一个自同态，h 是 E 的一个紧自同态。则 \(\mathrm{Sp}_{\omega}(u + h) = \mathrm{Sp}_{\omega}(u)\) ，且 \(\mathrm{Sp}_n(u + h) = \mathrm{Sp}_n(u)\) （对一切 \(n \in \overline{\mathbf{Z}} - \{0\}\) ）。

设 \(\lambda \in \mathbf{C}\) 。为使 \(u + h - \lambda 1_{\mathrm{E}}\) 是核（相应地，余核）有限维的严格态射，必须且只需 \(u - \lambda 1_{\mathrm{E}}\) 具有同样的性质，这由 III, p. 72 的定理 1（相应地，III, p. 73 的定理 2）得出。等式

\[
\begin{array}{c} \operatorname{Sp} _ {- \infty} (u + h) = \operatorname{Sp} _ {- \infty} (u), \quad \operatorname{Sp} _ {+ \infty} (u + h) = \operatorname{Sp} _ {+ \infty} (u), \\ \operatorname{Sp} _ {\omega} (u + h) = \operatorname{Sp} _ {\omega} (u) \end{array}
\]

由此得出。此外，我们有 \(\operatorname{Sp}_{n}(u+h)=\operatorname{Sp}_{n}(u)\) （由 III, p. 73 的定理 3，对每个 \(n\in \mathbf{Z}-\{0\}\) ）。

推论. --- 假设 \(\mathrm{Sp}_{\omega}(u)\) 的补集的无界连通分支包含 0。则 \(u + h\) 是 E 的里斯自同态。

有 \(\mathrm{Sp}_{\omega}(u+h)=\mathrm{Sp}_{\omega}(u)\) （由定理 2），故 0 属于 \(\mathbf{C}-\mathrm{Sp}_{\omega}(u+h)\) 的无界连通分支。由 III, p. 87 的定理 1 的推论，或者 0 不属于 \(u+h\) 的谱，或者它是该谱的敏感点。在两种情形下，\(u+h\) 都是 E 的里斯自同态。

\subsection{紧算子的谱}

引理 1. --- 设 E 是 R 上维数 \(\geqslant 2\) 的分离拓扑向量空间，X 是 E 的一个可数子集。补集 E - X 是连通的。

首先假设 E 的维数为 2。我们可以假设 \(E = R^{2}\) 赋有欧几里得范数（EVT, I, p. 14, 定理 2）。由于 X 可数，存在一个实数 \(r \in R_{+}^{*}\) ，使得以 0 为中心、r 为半径的圆 C 不与 X 相交。设 \(x \in E - X\) ；若 \(x \notin C\) ，存在一点 \(y \in C\) ，使得连接 x 与 y 的直线 \(L_{x}\) 不与 X 相交，因为 X 可数。集 E−X 是 C（它是连通的）与诸连通集 \(L_{x} \cup C\) （\(x \in \mathrm{E} - (\mathrm{X} \cup \mathrm{C})\) ）的并；这些集都包含 C，因此 E-X 是连通的（TG, I, p. 81, 命题 2）。

考察一般情形。对元素 \(x \in E - X\) 把 X 换成 X - x，我们化归到 \(0 \notin X\) 的情形。由于 E-X 是诸集 \(F - (X \cap F)\) （当 F 取遍 E 的 2 维子空间时所成的集）的并，且由前面的情形这些集都连通并包含 0，故集 E-X 是连通的（同前）。

引理 2. --- 设 S \(\subset \mathbf{C}\) 是 \(\mathbf{C} - \{0\}\) 中的一个无限的、离散的、有界的且闭的集。则 S 是某个序列 ( \(\lambda_{n}\) ) \(_{n\in N}\) 的值集，该序列由两两不同的非零复数组成，且序列 ( \(|\lambda_{n}|\) ) \(_{n\in N}\) 递减并收敛到 0。

对每个整数 \(i \geqslant 1\) ，集 \(\mathbf{A}_i\) （由复数 \(\lambda \in \mathrm{S}\) （满足 \(|\lambda| \geqslant \frac{1}{i}\) ）组成）在 \(\mathbf{C}\) 中是紧且离散的，故有限。用 \(a_i\) 表示它的基数。由于 S 无限，序列 \((a_i)\) 趋于 \(+\infty\) 。令 \(\mathbf{A}_0 = \varnothing\) ，\(a_0 = 0\) 。对每个 \(i \geqslant 1\) ，选取一个双射 \(n \mapsto \lambda_n\) （从区间 \([a_{i-1}, a_i[\) （\(\mathbf{N}\) 的）到 \(\mathbf{A}_i - \mathbf{A}_{i-1}\) 上），使得映射 \(n \mapsto |\lambda_n|\) 在 \([a_{i-1}, a_i[\) 上递减。序列 \((\lambda_n)_{n \in \mathbf{N}}\) 具有所要求的性质。

设 E 是无限维的完备可赋范空间。代数 \(\mathcal{L}^{\mathrm{c}}(\mathrm{E})\) 是 \(\mathcal{L}(\mathrm{E})\) 的一个不含单位元的子代数。回忆：对每个紧自同态 \(u \in \mathcal{L}^{\mathrm{c}}(\mathrm{E})\) ，谱 \(\mathrm{Sp}'_{\mathcal{L}^{\mathrm{c}}(\mathrm{E})}(u)\) 是 \(u\) 相对于含单位元子代数 \(\mathcal{L}^{\mathrm{c}}(\mathrm{E}) \oplus \mathbf{C1}_{\mathrm{E}}\) （\(\mathcal{L}(\mathrm{E})\) 的）的谱（I, p. 4, 第 4 目）。

命题 5. — 设 E 是完备可赋范空间，u 是 E 的一个紧自同态。\(\mathrm{Sp}_{\mathrm{s}}(u)\) 的每个元素都是有限谱重数的特征值。此外：

\begin{enumerate}
\def\labelenumi{\alph{enumi})}
\item
  若 E 有限维，则 \(\operatorname{Sp}(u) = \operatorname{Sp}_{\mathrm{s}}(u)\) ；
\item
  若 E 无限维，则 \(\mathrm{Sp}_{\mathrm{s}}(u) = \mathrm{Sp}(u) - \{0\}\) 且 \(\mathrm{Sp}_{\omega}(u) = \{0\}\) ；
\item
  若 \(\mathrm{Sp}_{\mathrm{s}}(u)\) 无限，则它是某个序列 \((\lambda_{n})_{n\in\mathbf{N}}\) 的值集，该序列由两两不同的非零复数组成，且序列 \((|\lambda_{n}|)_{n\in\mathbf{N}}\) 递减并收敛到 0；
\item
  若 E 无限维，则 \(\mathrm{Sp}(u) = \mathrm{Sp}'_{\mathcal{L}^c (\mathrm{E})}(u)\) 。
\end{enumerate}

\(\mathrm{Sp}_{\mathrm{s}}(u)\) 的每个元素都是有限谱重数的特征值（III, p. 83 的命题 1）。

断言 \(a\) 是初等的（III, p. 85, 例 1）。现在假设 E 无限维。

设 \(\lambda \in \mathrm{Sp}(u) - \{0\}\) 。则 \(u - \lambda 1_{\mathrm{E}}\) 是 E 的里斯自同态（III, p. 75 的命题 2 的推论 2），故 \(\lambda\) 属于 \(\mathrm{Sp}_{\mathrm{s}}(u)\) 。若 E 无限维，则 \(\mathrm{Sp}_{\omega}(u)\) 非空（III, p. 87, 定理 1, a)）。必然有 \(\mathrm{Sp}_{\omega}(u) = \{0\}\) ，由此得 \(b\) 。

集 \(\mathrm{Sp}_{\mathrm{s}}(u)\) 是离散且有界的。此外，我们有 \(\mathrm{Sp}_{\mathrm{s}}(u) = \mathrm{Sp}(u) \cap (\mathbf{C} - \{0\})\) （依照 \(b\) ），故 \(\mathrm{Sp}_{\mathrm{s}}(u)\) 在 \(\mathbf{C} - \{0\}\) 中是闭的。于是断言 \(c\) 由引理 2 得出。

\(u\) 的谱是可数的，故其在 \(\mathbf{C}\) 中的补集是连通的（引理 1）。把 I, p. 28 的命题 6 的推论应用于含单位元子代数 \(\mathcal{L}^{\mathrm{c}}(\mathrm{E})\oplus \mathbf{C}1_{\mathrm{E}}\) （\(\mathcal{L}(\mathrm{E})\) 的），我们断定 \(\mathrm{Sp}(u) = \mathrm{Sp}'_{\mathcal{L}^{\mathrm{c}}(\mathrm{E})}(u)\) 。

命题 6. --- 设 E 是 C 上的完备可赋范空间，u 是 E 的一个紧自同态。

\begin{enumerate}
\def\labelenumi{\alph{enumi})}
\item
  设 \(f \in \mathcal{O}(\mathrm{Sp}(u))\) 满足 \(f(0) = 0\) 。自同态 \(f(u)\) 是紧的；
\item
  假设 E 是复希尔伯特空间且 u 是正规的。设 f 是 \(\mathrm{Sp}(u)\) 上的一个连续函数，满足 \(f(0)=0\) 。正规自同态 \(f(u)\) 是紧的。
\end{enumerate}

此外，若 E 无限维，则逆命题成立；若 E 有限维，则条件 \(f(0) = 0\) 不是必要的。

我们可以假设 E 无限维。

我们来证明 \(a\) 及其逆。自同态 \(u\) 是巴拿赫代数 \(\mathcal{L}^{\mathrm{c}}(\mathrm{E})\) 的一个元素。由于 E 无限维，我们有等式 \(\mathrm{Sp}(u) = \mathrm{Sp}'_{\mathcal{L}^{\mathrm{c}}(\mathrm{E})}(u)\) （由命题 5 的 \(d\) ）。全纯泛函演算的元素 \(f(u)\) （该演算是把 \(\mathcal{L}^{\mathrm{c}}(\mathrm{E})\) 添加单位元后得到的含单位元巴拿赫代数的全纯泛函演算）属于 \(\mathcal{L}^{\mathrm{c}}(\mathrm{E})\) 当且仅当 \(f(0) = 0\) （I, p. 88）。而且这个元素与元素 \(f(u)\) （\(\mathcal{L}(\mathrm{E})\) 的）重合（I, p. 75 的命题 7），故 \(f(u)\) 是紧的当且仅当 \(f(0) = 0\) 。

断言 \(b\) 及其逆的证明完全相同，只需考察 C*-代数 \(\mathcal{L}^{\mathrm{c}}(\mathrm{E})\) 的连续泛函演算（I, p. 110, 定义 5）。

推论. --- 设 E 和 F 是希尔伯特空间，u 是从 E 到 F 中的一个连续线性映射。线性映射 u 是紧的，当且仅当 E 的自同态 \(|u|\) 是紧的。

设 \((j, |u|)\) 是 u 的极分解（I, p. 140 的定义 4）。由于 \(u = j|u|\) 且 \(|u| = \sqrt{u^{*}u}\) （I, p. 139 的命题 10），该等价性由 III, p. 5 的命题 3 以及前一命题的断言 b) 得出。

\subsection{希尔伯特空间的情形}

在本节中，E 表示 C 上的一个希尔伯特空间。我们用 \(\pi\) 表示从 \(\mathcal{L}(\mathrm{E})\) 到卡尔金代数 \(\mathcal{Calk}(\mathrm{E})\) 上的典范同态。

命题 7. --- 设 \(u \in \mathcal{L}(\mathrm{E})\) 。

\begin{enumerate}
\def\labelenumi{\alph{enumi})}
\item
  若 u 是正规的，则 u 是里斯自同态当且仅当 u 是指标为 0 的弗雷德霍姆自同态。
\item
  若 u 是埃尔米特的，则 u 是弗雷德霍姆自同态当且仅当 u 是里斯自同态。
\end{enumerate}

每个里斯自同态都是指标为 0 的弗雷德霍姆自同态（III, p. 46 的命题 5）。反之，假设 u 是指标为 0 的弗雷德霍姆自同态且 u 是正规的。此时它的幂零空间与它的核重合（EVT, V, p. 43, 命题 8），故有限维。于是由 III, p. 45 的引理 2 以及定义，u 是里斯自同态。

我们来证明 \(b\) 。由 \(a\) ，只需验证埃尔米特弗雷德霍姆自同态 \(u\) 的指标为零。但此时 \(u\) 的像（它是闭的）的正交补等于 \(u\) 的核（EVT, V, p. 41, 命题 4），由此得该断言。

推论. --- 设 \(u \in \mathcal{L}(\mathrm{E})\) 。若 u 是正规的，则 \(\mathrm{Sp}_{0}(u)\) 是空的；若 u 是埃尔米特的，则 \(\mathrm{Sp}_{\mathrm{e}}(u)\) 与 \(\pi(u)\) 相对于代数 \(\mathcal{Calk}(\mathrm{E})\) 的谱重合。

两个断言都由该命题以及诸集 \(\mathrm{Sp}_{n}(u)\) （\(n \in \overline{Z}\) ）与 \(\mathrm{Sp}_{\omega}(u)\) 的定义得出，它们构成 u 的本质谱的一个分划（III, p. 83 的定义 2）。

定理 3（外尔）. --- 设 \(u \in \mathcal{L}(\mathrm{E})\) 是 E 的一个正规自同态。\(u\) 的本质谱是诸集 \(\mathrm{Sp}(u + h)\) （当 \(h\) 取遍 \(\mathcal{L}^{\mathrm{c}}(\mathrm{E})\) 时）的交。

设 \(h \in \mathcal{L}^c(\mathrm{E})\) 。由于 \(\mathrm{Sp}_0(u)\) 是空的（上面的推论），III, p. 89 的定理 2 蕴含 \(\mathrm{Sp_e}(u + h) = \mathrm{Sp_e}(u)\) 。于是诸集 \(\mathrm{Sp}(u + h)\) 的交包含 \(\mathrm{Sp_e}(u)\) 。

设 \(\lambda \in \mathrm{Sp}_{\mathrm{s}}(u)\) 。用 \(\mathrm{E}_{\lambda}\) 表示 \(u\) 的相对于 \(\lambda\) 的特征空间，用 \(\mathrm{F}_{\lambda}\) 表示与 \(u\) 及 \(\mathbf{C} = \{\lambda\}\) 相伴的谱投影的像。空间 E 是 \(\mathrm{E}_{\lambda}\) 与 \(\mathrm{F}_{\lambda}\) 的拓扑直和。设 \(h\) 是 E 的这样一个有限秩自同态：它在 \(\mathrm{F}_{\lambda}\) 上为零，而在 \(\mathrm{E}_{\lambda}\) 上与恒同映射重合。自同态 \(u + h\) 可逆，故 \(\lambda \notin \mathrm{Sp}(u + h)\) 。定理由此得出。

于是，若 u 是 E 的正规自同态，且 \(h \in \mathcal{L}^{c}(\mathrm{E})\) ，则 \(u + h\) 的谱与 u 的谱只能相差一些具有有限谱重数的孤立点。